\documentclass[11pt]{article}
\usepackage[utf8]{inputenc}
\usepackage[T1]{fontenc}
\usepackage{amsmath,amssymb,amsthm,mathtools}
\usepackage[margin=1in]{geometry}
\usepackage{booktabs}
\usepackage{longtable}
\usepackage{array}
\newcolumntype{P}[1]{>{\raggedright\arraybackslash}p{#1}}
\usepackage{microtype}
\usepackage{enumitem}
\usepackage[font=small,labelfont=bf]{caption}
\usepackage{xcolor}
\usepackage{listings}
\usepackage{url}
\setlength{\emergencystretch}{2.5em}
\usepackage{hyperref}
\hypersetup{colorlinks=true, linkcolor=black, citecolor=black, urlcolor=blue}
\hypersetup{pdftitle={Every real primitive Dirichlet L-function has de Bruijn–Newman constant at most 9/32: a kernel-checked proof in Lean 4}, pdfauthor={Peter W. Murphy}}
\usepackage[capitalise,noabbrev]{cleveref}

\newtheorem{theorem}{Theorem}[section]
\newtheorem{proposition}[theorem]{Proposition}
\newtheorem{lemma}[theorem]{Lemma}
\newtheorem{corollary}[theorem]{Corollary}
\theoremstyle{definition}
\newtheorem{definition}[theorem]{Definition}
\newtheorem{remark}[theorem]{Remark}
\newtheorem{example}[theorem]{Example}


\newcommand{\lean}[1]{\texttt{#1}}
\newcommand{\re}{\operatorname{Re}}
\newcommand{\im}{\operatorname{Im}}
\newcommand{\C}{\mathbb{C}}
\newcommand{\R}{\mathbb{R}}
\newcommand{\N}{\mathbb{N}}
\newcommand{\Z}{\mathbb{Z}}
\newcommand{\GR}{\Gamma_{\mathbb{R}}}
\newcommand{\Lc}{\Lambda}
\newcommand{\todo}[1]{\textcolor{red}{[TODO: #1]}}

\definecolor{leanbg}{HTML}{F4F2EC}
\lstdefinelanguage{Lean}{
  morekeywords={theorem,lemma,def,noncomputable,import,namespace,end,open,in,by,fun,sorry,example,variable,scoped,Classical,structure,where},
  sensitive=true, morecomment=[l]{--}, morecomment=[s]{/-}{-/}, morestring=[b]",
}
\lstset{language=Lean, basicstyle=\small\ttfamily, backgroundcolor=\color{leanbg},
  frame=single, framerule=0pt, framesep=5pt, xleftmargin=5pt, xrightmargin=5pt,
  breaklines=true, columns=fullflexible, keepspaces=true, showstringspaces=false,
  literate={ℝ}{{$\mathbb{R}$}}1 {ℂ}{{$\mathbb{C}$}}1 {ℕ}{{$\mathbb{N}$}}1 {ℤ}{{$\mathbb{Z}$}}1
           {∀}{{$\forall$}}1 {∃}{{$\exists$}}1 {→}{{$\to$}}1 {≠}{{$\neq$}}1 {≤}{{$\le$}}1 {≥}{{$\ge$}}1
           {∧}{{$\wedge$}}1 {θ}{{$\theta$}}1 {Φ}{{$\Phi$}}1 {π}{{$\pi$}}1 {∑}{{$\sum$}}1 {∫}{{$\int$}}1
           {↦}{{$\mapsto$}}1 {←}{{$\leftarrow$}}1 {∈}{{$\in$}}1 {ξ}{{$\xi$}}1 {·}{{$\cdot$}}1
           {⟨}{{$\langle$}}1 {⟩}{{$\rangle$}}1 {₀}{{$_0$}}1 {²}{{$^2$}}1 {ε}{{$\varepsilon$}}1
           {σ}{{$\sigma$}}1 {τ}{{$\tau$}}1 {δ}{{$\delta$}}1 {γ}{{$\gamma$}}1 {ρ}{{$\rho$}}1
           {‖}{{$\|$}}1 {ψ}{{$\psi$}}1 {χ}{{$\chi$}}1 {↑}{{$\uparrow$}}1 {√}{{$\surd$}}1
           {⁻}{{$^{-}$}}1 {¹}{{$^{1}$}}1 {ℓ}{{$\ell$}}1 {∣}{{$\mid$}}1 {Λ}{{$\Lambda$}}1 {Γ}{{$\Gamma$}}1
           {𝓕}{{$\mathcal{F}$}}1 {𝓝}{{$\mathcal{N}$}}1 {¬}{{$\neg$}}1 {Σ}{{$\Sigma$}}1 {−}{{$-$}}1 {₄}{{$_4$}}1}

\title{Every real primitive Dirichlet $L$-function has\\ de Bruijn--Newman constant at most $9/32$:\\ a kernel-checked proof in Lean~4}
\author{Peter W. Murphy}
\date{Version 2.1.0, released 3 November 2026\\[2pt]
\small DOI of this version: \href{https://doi.org/10.5281/zenodo.23931764}{10.5281/zenodo.23931764}. Version 2.0.0: \href{https://doi.org/10.5281/zenodo.23858929}{10.5281/zenodo.23858929}; version 1.0.0: \href{https://doi.org/10.5281/zenodo.23840082}{10.5281/zenodo.23840082}; all versions: \href{https://doi.org/10.5281/zenodo.23840081}{10.5281/zenodo.23840081}\\[2pt]
\small Repository \url{https://github.com/DrMurphyIsIn/real-dirichlet-dbn-lean}\\[2pt]
\small The text of this paper is released under CC-BY-4.0; the Lean code under Apache-2.0.}

\begin{document}
\maketitle

\begin{abstract}
Let $\chi$ be a real primitive Dirichlet character of conductor $N \ge 3$ and parity $k \in \{0,1\}$, let $\Lc(s,\chi) = (N/\pi)^{(s+k)/2}\Gamma(\tfrac{s+k}{2})L(s,\chi)$ be its completed $L$-function, and let
\[
H^\chi_t(z) \;=\; \int_0^\infty e^{tu^2}\,\Phi_\chi(u)\cos(zu)\,du
\]
be the de Bruijn--Newman heat flow of $\Lc(\cdot,\chi)$, normalised so that $H^\chi_0(z) = \tfrac18\Lc(\tfrac12 + \tfrac{iz}{2},\chi)$. We give a proof, checked by the Lean~4 kernel, that for every such $\chi$ and every $t \ge 9/32$ all zeros of $H^\chi_t$ are real; equivalently, the generalised de Bruijn--Newman constant satisfies $\Lc_\chi \le 9/32$.

The mathematics is not new, and we say exactly whose it is. De Bruijn's strip-contraction theorem (1950) gives $\Lc_\chi \le D^2/2$ whenever the zeros of $H^\chi_0$ lie in $|\im z| \le D$. Freire, Gonz\'alez-Riquelme, Mu\~noz, Poblete and Salinas observed (arXiv:2610.27343, Remarks 3.1 and 3.4) that the zero-free half-plane $\re s > 7/8$ for Dirichlet $L$-functions asserted by OpenAI would give $D = 3/4$, and hence $\Lc_\chi \le 9/32$. What this paper adds is a complete formal proof of that composition. For conductors up to $400$, FGMPS's unconditional bounds are stronger for $136$ of the $242$ characters; what is new is that $9/32$, for every conductor, is a kernel-checked theorem. OpenAI's half-plane enters as their own kernel-checked Lean theorem. Every other step is proved in Lean on top of Mathlib: de Bruijn's theorem for a general class of kernels, the identification of the theta integral with Mathlib's completed $L$-function, and Gauss's sign theorem for real primitive characters. The final theorem depends only on the axioms \lean{propext}, \lean{Classical.choice} and \lean{Quot.sound}. In local, unsandboxed replays five checkers accept it (Lean's kernel, run twice, \lean{nanoda}, and the \lean{con-leche}/\lean{con-ron} pair), and all five reject a deliberately unsound control; and a sandboxed Comparator run against a separately written challenge accepts it (made by another agent session; the auditor reproduced its solution export byte for byte).

The result rests on OpenAI's Lean proof, whose accompanying manuscript is unrefereed. It is not progress on the Generalised Riemann Hypothesis, which is the statement $\Lc_\chi \le 0$. A short, separable section records two companion results for $\zeta$ that use no OpenAI code: $\Lc \le 225/512$ and $\psi(x) = x + O_\varepsilon(x^{31/32+\varepsilon})$, both from Cureton and Thorner's half-plane $\re s > 31/32$.

Version 2 adds three bounds on the same constant, in trust classes weaker than the kernel-checked theorem, and says which class each one is in (\cref{sec:v2}). (1) \emph{Unconditionally}, with no GRH, no zero-free half-plane and no numerical verification of zeros, and with a Siegel zero allowed if there is one: for every real primitive $\chi$ of conductor $N \ge 10^8$, $\Lc_\chi \le T(\ell) + \tfrac12 b(\ell) < 9/32$, where $\ell = \log(N/\pi)$ and $T$, $b$ are explicit elementary functions; the bound is $0.27643$ at $N = 10^8$ and is asymptotic to $\sqrt5\,\log\log N/\log N$. (2) Given OpenAI's half-plane, exactly as in version~1: $\Lc_\chi \le 0.27912 < 9/32$ for \emph{every} real primitive $\chi$, a strict improvement on $9/32$. (3) Unconditionally: $\Lc_\chi \le 0.49330$ for every real primitive $\chi$, which, as far as we could check, is the first bound below $1/2$ that is uniform in the conductor and uses no zero-free half-plane, no GRH and no zero verification. The method bounds the logarithmic derivative $-\im(H'_t/H_t)$ from below at a fixed height above the real axis by ball arithmetic (Arb), and feeds that bound into the zero-contraction inequalities of Planat and of Hakimoglu-Brown. These results are \emph{not} kernel-checked. They rest on Arb certificates that a separate AI auditor session reproduced (identical in every field except wall time), on AI-written proofs reviewed by one AI reviewer, and on unrefereed preprints. One step is kernel-checked: that the explicit function in (1) is below $9/32$ for every $N \ge 10^8$. Nothing in version~2 bears on GRH either.

Version 2.1.0 extends (1) downwards, by a sharper kernel in the same method (\cref{sec:W2B3}): unconditionally, $\Lc_\chi \le 0.2807226 < 9/32$ for every odd real primitive $\chi$ of conductor $N \ge 1{,}650{,}000$, and $\Lc_\chi \le 0.2808812 < 9/32$ for every even one of conductor $N \ge 3{,}700{,}000$. It applies Hakimoglu-Brown's derivative-kernel comparison (his Theorem~12) to the uniform floor, with a new Arb certifier for the vertical derivative of the logarithmic derivative, exact-rational Bernstein certificates of his polynomial positivity condition, and thirty-nine time slabs per parity. It is in the same trust classes as (1), its certificates were reproduced by the auditor, and nothing in version 2.0.0 is withdrawn. The one kernel-checked arithmetic statement still covers only $N \ge 10^8$; the middle range is not kernel-checked.
\end{abstract}

\tableofcontents

%======================================================================
\section{Introduction}\label{sec:intro}
%======================================================================

\subsection{The de Bruijn--Newman constant}

Let $\xi(s) = \tfrac12 s(s-1)\pi^{-s/2}\Gamma(s/2)\zeta(s)$ be Riemann's entire function. Following Rodgers and Tao \cite{RodgersTao2020} and Polymath15 \cite{Polymath2019}, put
\begin{equation}\label{eq:Hzeta}
H_t(z) = \int_0^\infty e^{tu^2}\,\Phi(u)\cos(zu)\,du,\qquad
\Phi(u) = \sum_{n\ge1}\bigl(2\pi^2n^4e^{9u} - 3\pi n^2e^{5u}\bigr)e^{-\pi n^2e^{4u}} .
\end{equation}
Then $H_0(z) = \tfrac18\xi(\tfrac12 + \tfrac{iz}{2})$, so the Riemann Hypothesis says that every zero of $H_0$ is real. The family $H_t$ solves the backward heat equation $\partial_t H_t = -\partial_z^2 H_t$. De Bruijn \cite{deBruijn1950} showed that real zeros stay real as $t$ increases and that every zero of $H_t$ is real for $t \ge 1/2$. Newman \cite{Newman1976} showed that some $H_t$ has non-real zeros. Together these give a constant $\Lc \in (-\infty, 1/2]$ such that every zero of $H_t$ is real exactly when $t \ge \Lc$. The Riemann Hypothesis is $\Lc \le 0$. Rodgers and Tao proved $\Lc \ge 0$ \cite{RodgersTao2020}. On the other side, Ki, Kim and Lee made de Bruijn's bound strict \cite{KKL2009}, Polymath15 proved $\Lc \le 0.22$ \cite{Polymath2019}, and Platt and Trudgian proved $\Lc \le 0.2$ \cite{PlattTrudgian2021}.

\subsection{The Dirichlet analogue}

For a real primitive character $\chi$ of conductor $N \ge 3$ the same construction applies to $\Lc(s,\chi)$. We follow the normalisation of Freire, Gonz\'alez-Riquelme, Mu\~noz, Poblete and Salinas \cite{FGMPS2026}, cited below as FGMPS. Their heat flow is $H^L_t(z) = \int_\R e^{tu^2+izu}\Phi_L(u)\,du$ with $H^L_0(z) = \xi_L(\tfrac{1+iz}{2})$, where $\xi_L$ is our $\Lc(\cdot,\chi)$ (\cite[(1.2), (1.3), (2.6)]{FGMPS2026}). Our $H^\chi_t$ equals $H^L_t/8$ (\cref{sec:normalisation}), so the two families have the same zeros at every time $t$.

FGMPS state that there is a finite constant $\Lc_\chi$ (their $\Lc_L$) such that every zero of $H^\chi_t$ is real if and only if $t \ge \Lc_\chi$; they credit the existence of $\Lc_\chi$ to Newman's work \cite{Newman1976}. The Generalised Riemann Hypothesis for $L(s,\chi)$ is equivalent to $\Lc_\chi \le 0$. Dobner proved $\Lc_\chi \ge 0$ for this family as part of a theorem for the whole extended Selberg class \cite{Dobner2021}, and de Bruijn's argument gives $\Lc_\chi \le 1/2$ from the critical strip. For the lower-bound side, FGMPS also survey work of Stopple \cite{Stopple2014}, of Andrade, Chang and Miller \cite{AndradeChangMiller2014}, and of Newman and Wu \cite{NewmanWu2020}; we have not read these papers and cite them only through FGMPS.

Before October 2026, the trivial bound $1/2$ was the only known upper bound \cite[\S1.4]{FGMPS2026}. FGMPS then proved the first non-trivial ones. They extend Polymath15's method to Dirichlet $L$-functions, using Platt's verification of GRH \cite{Platt2016} and interval arithmetic. They obtain the unconditional uniform bound
\[
\Lc_\chi \le 953/3200 = 0.2978125 \qquad (3 \le N \le 400),
\]
together with stronger individual bounds for each of the $242$ real primitive characters in that range \cite[Theorem 1.1, \S7.7]{FGMPS2026}. They also record a conditional statement, which is the one this paper formalizes. Remark~3.1 of \cite{FGMPS2026} reads:
\begin{quote}\small
``the zero-free half-plane $\re s > 7/8$ asserted in [24, Theorem 1.1], together with the functional equation, would imply the initial strip with $D = 3/4$. The conclusions reported for this value of $D$ are conditional on that global strip hypothesis.''
\end{quote}
Here [24] is OpenAI's preprint \cite{OAI7over8}. Remark~3.4 of \cite{FGMPS2026} adds: ``The initial strip hypothesis alone gives $\Lc_L \le D^2/2$ by Theorem 2.6''. Their Theorem 2.6 is de Bruijn's strip contraction, ``essentially contained in Theorem 13 of [9]'', where [9] is \cite{deBruijn1950}. At $D = 3/4$ this gives $\Lc_\chi \le 9/32$, for every conductor.

\subsection{OpenAI's half-plane}

On 30 September 2026 OpenAI released a preprint whose Theorem 1.1 asserts that every Dirichlet $L$-function, $\zeta$ included, has no zero in $\re s > 7/8$, apart from the pole of a principal character at $s = 1$ \cite{OAI7over8}. A companion preprint of 5 October, ``written with human assistance'', isolates the weaker half-plane $\re s > 11/12$ as an intermediate step and proves it separately \cite{OAIqrh}; FGMPS note that it would give $D = 5/6$ \cite[Remark~3.1]{FGMPS2026}. The release includes a Lean~4 library \cite{OAIlean} with the theorem
\begin{lstlisting}
theorem OAI.DirichletCharacter.LFunction_ne_zero_of_seven_eighths_lt_re
    {q : ℕ} [NeZero q] (χ : DirichletCharacter ℂ q) {s : ℂ}
    (hs : (7 / 8 : ℝ) < s.re) (hpole : ¬ (χ = 1 ∧ s = 1)) :
    _root_.DirichletCharacter.LFunction χ s ≠ 0
\end{lstlisting}
\noindent
The statement is copied from OpenAI's challenge file at commit \lean{adc7f124},
\[\text{\path{lean/ComparatorChallenges/DirichletSevenEighths.lean}},\]
where its proof is \lean{sorry}. The proof itself is in \path{lean/OAI/NumberTheory/DirichletL/}\allowbreak\path{Nonvanishing.lean}. The statement is about Mathlib's own \lean{DirichletCharacter.\allowbreak{}LFunction}, holds for every modulus and every character, and has no constants. Neither manuscript has been refereed. The Lean proof was replayed by the kernel in this programme on 7 October 2026; \cref{sec:trust} says exactly what that trust consists of.

\subsection{What this paper does}

\begin{theorem}[main theorem; Lean: \lean{bb\_\allowbreak{}real\_\allowbreak{}primitive\_\allowbreak{}meaning\_\allowbreak{}oai}]\label{thm:main}
Let $\chi$ be a primitive Dirichlet character modulo $N$, with $\chi \ne 1$ and $\chi$ real-valued, and let $k \in \{0,1\}$ be its parity. Then:
\begin{enumerate}[label=(\alph*)]
\item for every real $t \ge 9/32$, every zero of $H^\chi_t$ is real;
\item for every $z \in \C$, with $s = \tfrac12 + \tfrac{iz}{2}$,
\[
H^\chi_0(z) = \tfrac18\, N^{(s+k)/2}\,\Lc_{\mathrm{M}}(s,\chi) = \tfrac18\,\Lc(s,\chi),
\]
where $\Lc_{\mathrm{M}}(s,\chi)$ is Mathlib's \lean{DirichletCharacter.\allowbreak{}completedLFunction}, the entire function equal to $\GR(s+k)L(s,\chi)$ away from the poles of $\GR(s+k)$, and $\GR(s) = \pi^{-s/2}\Gamma(s/2)$.
\end{enumerate}
Consequently $\Lc_\chi \le 9/32$.
\end{theorem}

Part (a) is the bound. Part (b) is there so that a reader need check only the outer shape of our definition: that $H^\chi_t(z) = \int_0^\infty e^{tu^2}F(u)\cos(zu)\,du$ for one function $F$, independent of $t$ and $z$. Part (b) proves that at $t = 0$ this is $\tfrac18$ of the completed $L$-function of $\chi$, as Mathlib defines it. Since the identity holds for all $z \in \C$ (at $z = iy$ with $y < -1$ it forces $F\cosh(yu)$ to be integrable) and the cosine transform determines $F$, $F$ is the classical $\Phi_\chi$, whatever the formula in the definition says. \Cref{sec:statement} gives the Lean statement and a guide to reading it line by line.

\paragraph{What is ours, and what is not.} The theorem is FGMPS's conditional observation with its hypothesis discharged. The steps and their sources are:

\begin{center}\small
\begin{tabular}{@{}P{0.42\linewidth}P{0.25\linewidth}P{0.25\linewidth}@{}}
\toprule
Step & Mathematics due to & In Lean \\
\midrule
Zero-free half-plane $\re s > 7/8$ for $L(s,\chi)$ & OpenAI \cite{OAI7over8} & OpenAI \cite{OAIlean} \\
Half-plane plus functional equation give $D = 3/4$ & FGMPS, Rem.~3.1 \cite{FGMPS2026} (classical) & this work \\
Strip $D$ gives $\Lc_\chi \le D^2/2$ & de Bruijn \cite{deBruijn1950}; FGMPS Thm.~2.6, Rem.~3.4 & this work (generic) \\
Theta kernel, its Mellin transform, and the bridge to $\Lc(s,\chi)$ & classical & this work, on Mathlib's $L$-function API \\
Root number $\varepsilon(\chi) = 1$ for real primitive $\chi$ & Gauss \cite{Gauss1811} & this work \\
$L$-functions, Gauss sums, functional equation, Hurwitz kernels & classical & Mathlib \cite{Mathlib2020} \\
de Bruijn's theorem for $\zeta$, Hadamard products, Hurwitz's theorem & de Bruijn; classical & the \lean{dbn} island \cite{Murphy932} \\
\bottomrule
\end{tabular}
\end{center}

\noindent
The contribution of this paper is the kernel proof from end to end, together with three reusable formal components that, as far as we know, were not previously available in Lean:
\begin{enumerate}[label=(\roman*)]
\item de Bruijn's strip-contraction theorem for an abstract even kernel with double-exponential decay of arbitrary rate (\cref{sec:debruijn}). It generalises the \lean{dbn} island's theorem for $\zeta$ \cite{Murphy932}, with a Fourier-injectivity argument in place of the positivity of $\Phi$;
\item the bridge from the Dirichlet theta kernel to Mathlib's \lean{completedLFunction}, with the theta functional equation for real primitive characters (\cref{sec:theta});
\item Gauss's sign theorem, $\varepsilon(\chi) = 1$ for every real primitive $\chi$, in Mathlib's normalisation, including the closed form of the quadratic Gauss sum for every $N \ge 1$ (\cref{sec:gauss}). At Mathlib revision \lean{0493595e} there is no sign lemma for \lean{rootNumber} \cite[NOTE\_D2]{bbrecords}.
\end{enumerate}

\paragraph{What this paper does not do.} There is no new analytic input. The bound $9/32$ is the number FGMPS already wrote down conditionally. For $3 \le N \le 400$ it is not the best available: FGMPS's unconditional individual bounds beat $9/32$ for $136$ of the $242$ characters in that range \cite[\S7.7, Table 8]{FGMPS2026}. One example is $\chi_{-4}$, for which they prove $\Lc_{\chi_{-4}} \le 0.26$ \cite[\S7.5]{FGMPS2026}. Their bounds also do not depend on OpenAI. Our bound is uniform in $N$, but it is only as good as OpenAI's theorem, whose manuscript is unrefereed. Granting OpenAI's half-plane, FGMPS's own $D = 3/4$ search improves on $9/32$ for $149$ of the $242$ characters \cite[Table~8]{FGMPS2026}, and their Remark~3.4 gives $9/32$ for the rest and for every larger conductor. So no bound in this paper is better than what FGMPS's paper gives once OpenAI's theorem is granted; what is new is that the bound $9/32$, for every conductor, is a kernel-checked theorem. Nothing here bears on GRH, which is $\Lc_\chi \le 0$. We do not formalize Dobner's lower bound $\Lc_\chi \ge 0$, nor the existence of $\Lc_\chi$ as an infimum (Remark~\ref{rem:sinf}).

\subsection{What version 2 adds}\label{sec:v2intro}

Version 1 of this paper was a formal proof of a known conditional observation. Version 2 keeps that theorem and its verification record unchanged (\cref{sec:statement,sec:math,sec:formal,sec:verification} are as published, apart from cross-references), and adds three bounds on the same constant $\Lc_\chi$ that were obtained afterwards by a different method, with explicit dependence on the conductor. Version 2.1.0 adds a fourth, \cref{thm:Dintro}, which extends the first of them to smaller conductors. They are stated here informally; \cref{sec:v2} gives the exact statements, the method, the certificates and the records. Throughout, $\ell = \log(N/\pi)$.

\begin{theorem}[informal; \cref{thm:C2}]\label{thm:Aintro}
Let $\chi$ be a real primitive Dirichlet character of conductor $N \ge 10^8$. Then
\[
\Lc_\chi \;\le\; T(\ell) + \tfrac12\,b(\ell) \;<\; \tfrac{9}{32},
\]
where $T$ and $b$ are the explicit functions of \eqref{eq:C2bound}. Nothing is assumed beyond $\chi$ being real and primitive: no GRH, no zero-free half-plane, no numerical verification of zeros; a Siegel zero, if there is one, is allowed. At $N = 10^8$ the bound is $0.27643$; at $\ell = 10^3$ it is $0.0123$; asymptotically it is $(\sqrt5 + o(1))\log\log N/\log N$.
\end{theorem}

\begin{theorem}[informal; \cref{thm:Y3}]\label{thm:Bintro}
Assume, exactly as in \cref{thm:main}, OpenAI's zero-free half-plane $\re s > 7/8$ for $L(s,\chi)$. Then $\Lc_\chi \le 0.27912 < 9/32$ for every real primitive $\chi \ne 1$: at most $0.2791114$ for odd $\chi$ and $0.2783189$ for even $\chi$ when $N < 10^8$, and at most $0.27643$ when $N \ge 10^8$, by \cref{thm:Aintro}.
\end{theorem}

\begin{theorem}[informal; \cref{thm:U1}]\label{thm:Cintro}
Unconditionally, $\Lc_\chi \le 0.49330$ for every real primitive $\chi \ne 1$.
\end{theorem}

\begin{theorem}[informal; \cref{thm:W2B3}; new in version 2.1.0]\label{thm:Dintro}
Unconditionally, with the same hypotheses as \cref{thm:Aintro}: $\Lc_\chi \le 0.2807226 < 9/32$ for every odd real primitive $\chi$ of conductor $N \ge 1{,}650{,}000$, and $\Lc_\chi \le 0.2808812 < 9/32$ for every even one of conductor $N \ge 3{,}700{,}000$. For $N \ge 10^8$, \cref{thm:Aintro} is the better bound.
\end{theorem}

\noindent
The method is not the one of \cref{sec:math}. It lower-bounds $S_\chi(x,p,t) = -\im(H'_t/H_t)(x+ip)$, the logarithmic derivative of the heat flow at a fixed height $p$ above the real axis, uniformly in $x$ and over a range of $t$, by interval (ball) arithmetic, and feeds the bound into the contraction inequalities for the highest non-real zero due to Planat \cite{Planat2026} and Hakimoglu-Brown \cite{HakimogluBrown2026}. Where \cref{thm:main} used only the functional equation and complex analysis, these results use real computation. The conductor enters $S_\chi$ explicitly only through the term $\ell/4$, and otherwise only through the centre of a Gaussian average, which $\ell$ moves to the right; that is why the bounds improve with $N$ and why, from $N = 10^8$ on with one contraction step (\cref{thm:Aintro}), and from $N = 1{,}650{,}000$ (odd) or $3{,}700{,}000$ (even) on with thirty-nine steps and the derivative kernel (\cref{thm:Dintro}), no hypothesis on the zeros is needed at all.

\paragraph{Trust classes.} This paper now contains results in four distinct trust classes, and we keep them apart throughout.
\begin{enumerate}[label=(T\arabic*)]
\item\label{T1} \emph{Kernel-checked.} \Cref{thm:main}, as in version~1: a Lean~4 theorem replayed by five checkers and accepted by a sandboxed Comparator run. Also, new in version~2, the single arithmetic statement that the explicit function of \cref{thm:Aintro} is below $9/32$ for every $N \ge 10^8$ (\lean{C2.C2bound\_\allowbreak{}lt\_\allowbreak{}9\_\allowbreak{}32}, \cref{sec:v2lean}). That statement was checked by the Lean kernel in one build, in a different environment from version~1's, and was not replayed by further checkers. Version 2.1.0 adds nothing to this class: the arithmetic statement covers $N \ge 10^8$ only, and the middle range of \cref{thm:Dintro} is not kernel-checked at all.
\item\label{T2} \emph{Arb-certified.} Lower bounds for $S_\chi$ over boxes of parameters, computed in ball arithmetic with outward rounding \cite{flint}, with every box, its result and the failure count recorded in a ledger. The auditor re-ran every ledger that \cref{sec:v2} depends on and obtained ledgers identical in every field except wall time (REVIEW\_Y3, REVIEW\_Y3\_G1 and, for the two P5 components, REVIEW\_PRERELEASE\_v2.0.0 \cite{dq1records}); for version 2.1.0, the seventy-eight slab ledgers of \cref{sec:W2B3}, which also carry upper bounds for $\partial_pS_\chi$, were re-run by the auditor field by field (REVIEW\_W2B3), and the exact-rational polynomial certificates there were reproduced by the auditor's own checker. These are rigorous numerical certificates, not kernel-checked proofs: their soundness depends on the certifier programs, which are AI-written and were reviewed by one AI reviewer; four defects in an earlier version of them were found, three by a second agent session and one by the reviewer (\cref{sec:defects}), and a fifth, in the new certifier of version 2.1.0, by the reviewer (\cref{sec:W2B3}).
\item\label{T3} \emph{AI-written proofs, reviewed by one AI reviewer.} The analytic steps that connect the certificates to the bounds: the Gaussian-tilt identity for $S_\chi$, the reduction of the domain, the tail lemmas, the confinement lemma, and the assembly with the cited contraction theorems; for version 2.1.0, the variance form of $\partial_pS_\chi$, its tail bookkeeping, and the transposition of Hakimoglu-Brown's stepped derivative-kernel comparison (\cref{sec:W2B3}). No human has refereed them.
\item\label{T4} \emph{Cited.} De Bruijn's strip theorem; FGMPS's heat-kernel form, finiteness theorem and effective approximation; Planat's and Hakimoglu-Brown's contraction inequalities, and, for \cref{thm:Dintro}, Hakimoglu-Brown's conjugate-pair inequality with a derivative term (his Lemma~10) and stepped comparison (his Theorem~12); Euler-product and convexity bounds for $L(s,\chi)$, read in Bennett, Martin, O'Bryant and Rechnitzer \cite{BMOR2020}; three formulas from the DLMF \cite{DLMF}; and OpenAI's half-plane, for \cref{thm:Bintro} only. FGMPS, Planat, Hakimoglu-Brown and OpenAI are unrefereed preprints (Hakimoglu-Brown's is a working draft), and the two contraction papers are written for Riemann's $\xi$; their transposition to Dirichlet $L$-functions is in class \ref{T3}.
\end{enumerate}
\Cref{sec:v2trust} assigns each step of the new results to its class. Nothing in version~2 is RH or GRH progress; every statement is a bound on $\Lc_\chi$ strictly above~$0$.

\paragraph{Version history.} Version 1.0.0 (31 October 2026) is \cref{thm:main} with its verification record. Version 2.0.0 (1 November 2026) added \cref{sec:v2} with \cref{thm:C2,thm:Y3,thm:U1}, the certificates and the Lean arithmetic step, and changed nothing in version~1. Version 2.1.0 adds \cref{thm:W2B3} (the middle range of conductors, \cref{sec:W2B3}) with its certificates in \path{certificates/w2b3/}, and the corresponding sentences in the abstract, this introduction, the trust classes, Remark~\ref{rem:U2}, \cref{sec:v2repro,sec:limits,sec:ai} and Appendix~\ref{app:v2records}. Nothing in version 2.0.0 is withdrawn: every theorem, number and record of that version stands as published, and Remark~\ref{rem:U2} is improved, not corrected.

\paragraph{How to read this paper.} \Cref{sec:statement} states the kernel-checked theorem formally and explains how to read it. \Cref{sec:math} gives its mathematics in full, for a number theorist who does not read Lean. \Cref{sec:formal} describes the formalization and its trust base, \cref{sec:verification} the checks it passed, and \cref{sec:companion} the optional companion results for $\zeta$. \Cref{sec:v2} is version~2: the method, the four new bounds, their certificates, the defects found and repaired on the way, and the one arithmetic step that is kernel-checked. \Cref{sec:limits} lists the limitations and \cref{sec:ai} states how AI was used.

%======================================================================
\section{The statement}\label{sec:statement}
%======================================================================

\subsection{The definitions}

Four non-Mathlib definitions appear in the statement. Here they are, excerpted verbatim from the challenge file \lean{ChallengeDefs.lean} (md5 \lean{67e538bc}; only its \lean{namespace} and \lean{open} lines and the docstring of \lean{Hchi} are omitted), which copies them from the audited sources \lean{GenTheta.lean} (md5 \lean{3dbf6dec}) and \lean{GenThetaFE.lean} (md5 \lean{8aeec9b0}):

\begin{lstlisting}
variable (c : ℕ → ℝ) (k : ℕ) (q : ℝ)

/-- `θ_c(x) = Σ_n c(n) n^k exp(−π n² x / q)`. -/
noncomputable def thetaC (x : ℝ) : ℝ :=
  ∑' n : ℕ, c n * (n : ℝ) ^ k * Real.exp (-π * (n : ℝ) ^ 2 * x / q)

/-- `Φ_c(u) = e^{(1+2k)u} θ_c(e^{4u})`. -/
noncomputable def PhiC (u : ℝ) : ℝ :=
  Real.exp ((1 + 2 * k) * u) * thetaC c k q (Real.exp (4 * u))

variable {N : ℕ} [NeZero N] (χ : DirichletCharacter ℂ N)

variable {χ} {c : ℕ → ℝ} {k : ℕ}

open scoped Classical in
/-- The parity weight `k(χ) = 0` (even) / `1` (odd). -/
noncomputable def parityK (χ : DirichletCharacter ℂ N) : ℕ := if χ.Even then 0 else 1

open scoped Classical in
noncomputable def Hchi (χ : DirichletCharacter ℂ N) (t : ℝ) (z : ℂ) : ℂ :=
  ∫ u in Ioi (0 : ℝ), ((Real.exp (t * u ^ 2) : ℝ) : ℂ) *
    ((PhiC (fun n ↦ (χ n).re) (parityK χ) N u : ℝ) : ℂ) * Complex.cos (z * u)
\end{lstlisting}

\noindent
In mathematical notation, with $c(n) = \re\chi(n)$, $k = k(\chi)$ and $q = N$:
\begin{equation}\label{eq:defs}
\begin{gathered}
\theta_\chi(x) = \sum_{n\ge0}\chi(n)\,n^k e^{-\pi n^2x/N},\qquad
\Phi_\chi(u) = e^{(1+2k)u}\,\theta_\chi(e^{4u}),\\
H^\chi_t(z) = \int_0^\infty e^{tu^2}\Phi_\chi(u)\cos(zu)\,du .
\end{gathered}
\end{equation}
Three points about these transcriptions:
\begin{itemize}
\item The sum is a Mathlib \lean{tsum} over all $n \in \N$, including $n = 0$. The $n = 0$ term vanishes because $\chi(0) = 0$ when $N \ne 1$, and $\chi \ne 1$ forces $N \ne 1$. The proofs derive this internally.
\item For real-valued $\chi$, $c(n) = \re\chi(n) = \chi(n)$. The definition takes the real part so that $\Phi_\chi$ is real-valued by construction. It makes sense for every Dirichlet character, but it is only meaningful for real ones.
\item The integral is a Bochner integral over $(0,\infty)$ for Lebesgue measure. \lean{Complex.cos} is the entire cosine, so $H^\chi_t(z)$ is defined for complex $z$. Integrability for every real $t$ and complex $z$ is proved, from the decay estimate \eqref{eq:decay} below. A Bochner integral of a non-integrable function would silently be $0$, so this proof matters.
\end{itemize}

\subsection{The theorem}

The theorem, as stated and proved in OpenAI's environment (file \lean{BbOaiMeaning.lean}, md5 \lean{23ea25e9}), is:

\begin{lstlisting}[numbers=left,numberstyle=\tiny,numbersep=4pt]
open Complex in
theorem bb_real_primitive_meaning_oai :
    (∀ {N : ℕ} [NeZero N] (χ : DirichletCharacter ℂ N), χ.IsPrimitive → χ ≠ 1 → (∀ n : ℕ, (χ n).im = 0) →
      ∀ t : ℝ, 9 / 32 ≤ t → ∀ z : ℂ, DBN.Gen.Hchi χ t z = 0 → z.im = 0) ∧
    (∀ {N : ℕ} [NeZero N] (χ : DirichletCharacter ℂ N), χ.IsPrimitive → χ ≠ 1 → (∀ n : ℕ, (χ n).im = 0) → ∀ z : ℂ,
      DBN.Gen.Hchi χ 0 z = (1 / 8) * (N : ℂ) ^ ((1 / 2 + I * z / 2 + DBN.Gen.parityK χ) / 2) *
        DirichletCharacter.completedLFunction χ (1 / 2 + I * z / 2)) :=
  ⟨bb_real_primitive_real_zeros_oai, fun χ hprim hχ1 hreal z ↦ DBN.Gen.Hchi_bridge χ hprim hχ1 hreal z⟩
\end{lstlisting}

\paragraph{Reading guide.}
\begin{center}\small
\begin{longtable}{@{}p{0.08\linewidth}p{0.86\linewidth}@{}}
\toprule
Line & Reading \\
\midrule
1 & Inside this declaration, \lean{I} means the imaginary unit \lean{Complex.I}. \\
2 & The name of the theorem. The prefix \lean{bb\_\allowbreak{}} records that the auditor wrote this composition file. \\
3 & For every modulus $N \ge 1$ (\lean{NeZero N} says $N \ne 0$; the braces make $N$ implicit) and every Dirichlet character $\chi$ modulo $N$ with values in $\C$ that is primitive in Mathlib's sense (\lean{IsPrimitive}: its conductor is $N$), is not the trivial character, and takes only real values at the natural numbers, \dots \\
4 & \dots and for every real $t \ge 9/32$ and every complex $z$ with $H^\chi_t(z) = 0$, we have $\im z = 0$. This is part (a) of \cref{thm:main}. The numerals \lean{9 /\allowbreak{} 32} are elaborated in $\R$; \lean{\#check} with \lean{pp.numericTypes} prints \lstinline[basicstyle=\ttfamily,backgroundcolor={}]|(9 / 32 : ℝ)|. \\
5 & And, with the same hypotheses on $\chi$, for every complex $z$: \\
6--7 & $H^\chi_0(z) = \tfrac18 \cdot N^{(s+k)/2}\cdot\Lc_{\mathrm{M}}(s,\chi)$, where $s = \tfrac12 + \tfrac{iz}{2}$, $k = k(\chi)$, and $N^{(\cdot)}$ is Mathlib's complex power with base the natural number $N$ cast to $\C$. Here \lean{completedLFunction} is Mathlib's, from \lean{Mathlib.NumberTheory.LSeries.DirichletContinuation}. This is part (b). \\
8 & The proof: the pair of the two theorems \lean{bb\_\allowbreak{}real\_\allowbreak{}primitive\_\allowbreak{}real\_\allowbreak{}zeros\_\allowbreak{}oai} and \lean{Hchi\_\allowbreak{}bridge}, described in \cref{sec:formal}. \\
\bottomrule
\end{longtable}
\end{center}

\noindent
Part (a) alone uses $43$ constants ($29$ from Mathlib, $13$ from \lean{Init}, and \lean{Hchi}); the whole conjunction uses $59$ ($37$ from Mathlib, $20$ from \lean{Init}, and \lean{Hchi} and \lean{parityK} from \lean{GenThetaFE}). Every one of them comes from Lean's core library (\lean{Init}), from Mathlib, or from \lean{GenThetaFE} (\lean{Hchi} and \lean{parityK}). None comes from OpenAI's \lean{OAI.*} namespace. The auditor's program, which lists each constant of the statement together with the module that declares it, confirms this \cite[REVIEW\_GAUSS\_SIGN\_AND\_OAI \S B.3; ADD1 \S3]{bbrecords}.

\subsection{Mathlib's completed $L$-function, and the classical one}\label{sec:mathlibLambda}

Mathlib's $\Lc_{\mathrm{M}}(s,\chi) = $ \lstinline[basicstyle=\ttfamily,backgroundcolor={}]|completedLFunction χ s| omits the conductor factor. Its relevant facts, read from Mathlib's source at revision \lean{0493595e} \cite[ADD1 \S5]{bbrecords}, are:
\begin{itemize}
\item \lstinline[basicstyle=\ttfamily,backgroundcolor={}]|gammaFactor χ s = if χ.Even then Gammaℝ s else Gammaℝ (s + 1)|, where $\GR(s) = \pi^{-s/2}\Gamma(s/2)$;
\item \lean{LFunction\_\allowbreak{}eq\_\allowbreak{}completed\_\allowbreak{}div\_\allowbreak{}gammaFactor}: $L(s,\chi) = \Lc_{\mathrm{M}}(s,\chi)/\GR(s+k)$ when $s \ne 0$ or $N \ne 1$;
\item \lean{IsPrimitive.completedLFunction\_\allowbreak{}one\_\allowbreak{}sub}: for primitive $\chi$,
\begin{equation}\label{eq:FE}
\Lc_{\mathrm{M}}(1-s,\chi) = N^{s-1/2}\,\varepsilon(\chi)\,\Lc_{\mathrm{M}}(s,\chi^{-1}),
\end{equation}
with \lstinline[basicstyle=\ttfamily,backgroundcolor={}]|rootNumber χ| $= \varepsilon(\chi) = \tau(\chi)/(i^k N^{1/2})$, where $\tau(\chi) = \sum_{a \bmod N}\chi(a)e^{2\pi i a/N}$ (\lstinline[basicstyle=\ttfamily,backgroundcolor={}]|gaussSum χ stdAddChar|).
\end{itemize}
Hence $N^{(s+k)/2}\Lc_{\mathrm{M}}(s,\chi) = (N/\pi)^{(s+k)/2}\Gamma(\tfrac{s+k}2)L(s,\chi) = \Lc(s,\chi)$ is the classical completed $L$-function. It satisfies $\Lc(1-s,\chi) = \varepsilon(\chi)\Lc(s,\bar\chi)$, because the powers of $N$ in \eqref{eq:FE} combine. Since $N^{(s+k)/2} \ne 0$, part (b) says that the zeros of $H^\chi_0$ are exactly the points $z$ with $\Lc(\tfrac12 + \tfrac{iz}2,\chi) = 0$. These points are in bijection with the non-trivial zeros of $L(s,\chi)$. The identification of $N^{(s+k)/2}\Lc_{\mathrm{M}}$ with the classical $\Lc$ is this reading of Mathlib's definitions; the Lean statement itself is about Mathlib's constant.

\subsection{Relation to the FGMPS normalisation}\label{sec:normalisation}

In \cite[(1.3)]{FGMPS2026}, $\Phi_L(u) = 4\varphi_L(e^{4u})$ with $\varphi_L(x) = x^{1/4+a/2}\psi_L(x)$ and $\psi_L(x) = \sum_{n\ge1}n^a\chi(n)e^{-n^2\pi x/m}$. Here $m = N$ and $a = k$, and $x^{1/4 + a/2}$ at $x = e^{4u}$ is $e^{(1+2k)u}$, so $\Phi_L = 4\Phi_\chi$. For even $\Phi_\chi$ (\cref{sec:theta}), $H^L_t(z) = \int_\R e^{tu^2}\Phi_L(u)e^{izu}\,du = 8H^\chi_t(z)$ for every real $t$. So both families have the same zeros at every time, and the constant $\Lc_\chi$ is the same in both conventions. This identity for $t \ne 0$ is a one-line unfolding of the definitions, not a Lean theorem: FGMPS's $H^L_t$ is not defined in Lean. What is formal is part (b), which anchors $\Phi_\chi$ to $\Lc(s,\chi)$.

\begin{remark}[the infimum-free form]\label{rem:sinf}
As in our earlier papers \cite{Murphy932,MurphyNewman}, ``$\Lc_\chi \le 9/32$'' is stated as ``for every $t \ge 9/32$ every zero of $H^\chi_t$ is real''. Once $\Lc_\chi$ exists and de Bruijn's monotonicity holds, the two are equivalent. Monotonicity is in the formalization (\lean{dbn\_\allowbreak{}real\_\allowbreak{}zeros\_\allowbreak{}upset}); it is stated for abstract kernels and transfers to $H^\chi_t$ through the rescaling of \cref{sec:debruijn}. Existence, which rests on Newman's theorem, is not. Stating the bound without an infimum keeps the theorem a statement about the zeros of an explicit entire function.
\end{remark}

%======================================================================
\section{The mathematics}\label{sec:math}
%======================================================================

This section is self-contained. Every theorem, proposition and corollary in it is proved in the formalization; the two places where a remark is not formalized say so. Lean names are given in brackets for readers who want to find them.

\subsection{The theta kernel and the bridge}\label{sec:theta}

Throughout, $\chi$ is a non-trivial real Dirichlet character modulo $N$ with parity $k$: $\chi(-1) = (-1)^k$, so $\Lc_{\mathrm{M}}(s,\chi) = \GR(s+k)L(s,\chi)$ away from the poles of $\GR(s+k)$.

\paragraph{Decay.} Since $|\chi(n)| \le 1$ and $\chi(0) = 0$, for $x \ge 1$
\[
|\theta_\chi(x)| \le \sum_{n\ge1}n^k e^{-\pi n^2x/N} \le e^{-\pi x/(2N)}\sum_{n\ge1}n^ke^{-\pi n^2/(2N)},
\]
using $n^2x \ge (n^2+x)/2$ for $n, x \ge 1$. Hence, for $u \ge 0$,
\begin{equation}\label{eq:decay}
|\Phi_\chi(u)| \le S_{k,N}\; e^{(1+2k)u}\exp\!\bigl(-\tfrac{\pi}{2N}e^{4u}\bigr)
\qquad\text{[\lean{abs\_\allowbreak{}PhiC\_\allowbreak{}le}]}.
\end{equation}
The true rate is $\pi/N$; the weaker rate suffices for everything below. $\Phi_\chi$ is continuous [\lean{continuous\_\allowbreak{}PhiC}]. It is not identically zero, because the $n = 1$ term dominates as $u \to \infty$ [\lean{exists\_\allowbreak{}PhiC\_\allowbreak{}ne\_\allowbreak{}zero}].

\paragraph{The Mellin transform.} For $\re s > 1$, integrating term by term,
\[
\int_0^\infty \theta_\chi(x)\,x^{(s+k)/2}\,\frac{dx}{x}
= \sum_{n\ge1}\chi(n)n^k\Bigl(\frac{N}{\pi n^2}\Bigr)^{(s+k)/2}\Gamma\Bigl(\frac{s+k}2\Bigr)
= \Lc(s,\chi)
\]
[\lean{mellin\_\allowbreak{}thetaC}, \lean{div\_\allowbreak{}pi\_\allowbreak{}cpow\_\allowbreak{}mul\_\allowbreak{}Gamma}].
Put $x = e^{4u}$ and $s = \tfrac12 + \tfrac{iz}2$. Then $x^{(s+k)/2} = e^{(1+2k)u}e^{izu}$ and $dx/x = 4\,du$, so
\begin{equation}\label{eq:fourier}
\Lc(\tfrac12 + \tfrac{iz}2,\chi) = 4\int_\R \Phi_\chi(u)\,e^{izu}\,du\qquad(\im z < -1)\qquad\text{[\lean{fourier\_\allowbreak{}PhiC\_\allowbreak{}eq\_\allowbreak{}mellin}]}.
\end{equation}

\paragraph{Evenness and the bridge.} If $\Phi_\chi$ is even, the two-sided integral in \eqref{eq:fourier} is twice the one-sided cosine integral [\lean{Hd\_\allowbreak{}zero\_\allowbreak{}eq\_\allowbreak{}half\_\allowbreak{}fourier}]. So $H^\chi_0(z) = \tfrac18\Lc(\tfrac12 + \tfrac{iz}{2},\chi)$ for $\im z < -1$. Both sides are entire: the left by differentiation under the integral, using \eqref{eq:decay}, and the right because $\chi \ne 1$. Since they agree on an open set, they agree everywhere [\lean{DBN.Gen.bridge}]. Primitivity is not used here, and evenness is used only to fold the integral. The formula holds for every real $\chi \ne 1$ whose kernel is even.

\paragraph{The theta functional equation.} For real primitive $\chi$, Poisson summation gives
\begin{equation}\label{eq:thetaFE}
\theta_\chi(1/x) = \varepsilon(\chi)\,x^{(1+2k)/2}\,\theta_\chi(x)\qquad(x > 0).
\end{equation}
The formalization proves this with $\varepsilon(\chi) = 1$ as a hypothesis [\lean{theta\_\allowbreak{}fe}]. It splits $\sum_{n\in\Z}$ into residue classes modulo $N$, applies Mathlib's functional equations for the even and odd Hurwitz theta kernels, and recombines the classes with Mathlib's Gauss-sum identity for the Fourier transform of a primitive character (\lean{IsPrimitive.fourierTransform\_\allowbreak{}eq\_\allowbreak{}inv\_\allowbreak{}mul\_\allowbreak{}gaussSum}), using $\chi^{-1} = \chi$. Substituting $x = e^{4u}$ into \eqref{eq:thetaFE} gives
\[
\Phi_\chi(-u) = \varepsilon(\chi)\,\Phi_\chi(u)\qquad\text{[\lean{PhiC\_\allowbreak{}even\_\allowbreak{}of\_\allowbreak{}rootNumber}, for } \varepsilon = 1\text{]}.
\]
So $\Phi_\chi$ is even exactly when $\varepsilon(\chi) = 1$. Gauss's sign theorem supplies this.

\subsection{Gauss's sign theorem}\label{sec:gauss}

\begin{theorem}[Gauss; Lean: \lean{DBN.Gen.rootNumber\_\allowbreak{}eq\_\allowbreak{}one\_\allowbreak{}of\_\allowbreak{}real\_\allowbreak{}primitive}]\label{thm:gauss}
Every real primitive Dirichlet character $\chi$ has $\varepsilon(\chi) = 1$, that is, $\tau(\chi) = \sqrt N$ if $\chi$ is even and $\tau(\chi) = i\sqrt N$ if $\chi$ is odd.
\end{theorem}

\begin{lstlisting}
theorem DBN.Gen.rootNumber_eq_one_of_real_primitive {N : ℕ} [NeZero N] (χ : DirichletCharacter ℂ N)
    (hprim : χ.IsPrimitive) (hreal : ∀ n : ℕ, (χ n).im = 0) :
    DirichletCharacter.rootNumber χ = 1
\end{lstlisting}

\noindent
The statement uses only constants from Mathlib and \lean{Init}; \lean{rootNumber} is Mathlib's own definition \cite[REVIEW\_GAUSS\_SIGN\_AND\_OAI \S A.3]{bbrecords}. Gauss found the sign of the quadratic Gauss sum in 1805 (his letter to Olbers of 3 September 1805 is the usual reference); the proof, \emph{Summatio quarumdam serierum singularium}, was presented in 1808 and printed in 1811 \cite{Gauss1811}. We have not checked these dates against the originals; for the history of the theorem see Berndt, Evans and Williams \cite{BEW1998}. The classical deduction for real primitive characters goes through their identification with Kronecker symbols of fundamental discriminants; see \cite[Theorems 9.13, 9.16]{MontgomeryVaughan2006}, as cited in \cite[Lemma 2.9]{FGMPS2026}. The formalization follows a different, self-contained route in three stages, 1,419 lines in all:
\begin{enumerate}[label=S\arabic*.]
\item \emph{The quadratic Gauss sum.} For every $N \ge 1$, $G(N) = \sum_{j=0}^{N-1}e^{2\pi ij^2/N} = \tfrac12(1+i)(1+(-i)^N)\sqrt N$ [\lean{gaussG\_\allowbreak{}eq}]. The proof is analytic. It evaluates the Jacobi theta function at $\tau_t = 2/N + it$ as $t \to 0^+$ in two ways: directly, by residue classes, and through the transformation $\tau \mapsto -1/\tau$.
\item \emph{Prime conductor.} For an odd prime $p$, a real character $\chi \ne 1$ modulo $p$ is the Legendre symbol and is primitive. Its Gauss sum is $G(p)$, which gives $\varepsilon = 1$ [\lean{rootNumber\_\allowbreak{}eq\_\allowbreak{}one\_\allowbreak{}of\_\allowbreak{}prime}].
\item \emph{Induction on the conductor.} Strong induction on $N$ peels off one odd prime $p \parallel N$ at a time, writing $\chi = \chi_p\chi'$ by the Chinese remainder theorem. Alongside $\varepsilon = 1$ it carries the invariant $\chi(\ell) = \bigl(\tfrac{\chi(-1)N}{\ell}\bigr)$ for odd primes $\ell \nmid N$. With this invariant, the cross term $\chi_p(N/p)\chi'(p)$ reduces to $\bigl(\tfrac{\chi'(-1)}{p}\bigr)$ without a further appeal to reciprocity; quadratic reciprocity is used only to establish the invariant for the Legendre factor. The base cases are the levels $1$, $4$ and $8$, that is, the characters $1, \chi_{-4}, \chi_8, \chi_{-8}$. The levels $2$ and $2^e$ with $e \ge 4$, and any level divisible by $p^2$ for an odd prime $p$, carry no real primitive character [\lean{gauss\_\allowbreak{}sign\_\allowbreak{}aux}].
\end{enumerate}
The theorem quantifies over Mathlib's \lean{IsPrimitive} directly. It does not use the correspondence between real primitive characters and fundamental discriminants.

\begin{remark}[the sign is needed for evenness, not for the strip]
The zero strip of \cref{sec:strip} does not depend on $\varepsilon$, because \eqref{eq:FE}, with $\chi^{-1} = \chi$, sends a zero at $s$ to a zero at $1 - s$ whatever the value of $\varepsilon$. The sign matters only because the formalized de Bruijn theorem is stated for even kernels. The auditor noted an alternative \cite[NOTE\_D2 \S2]{bbrecords}. For real primitive $\chi$, applying \eqref{eq:FE} twice gives $\varepsilon^2 = 1$ without any Gauss-sum work, and the case $\varepsilon = -1$ would need a version of de Bruijn's theorem for odd kernels ($z\,K(z^2)$ in place of $K(z^2)$). This route was not formalized (it is one of the two unformalized remarks of this section); Gauss's theorem was proved instead.
\end{remark}

\subsection{De Bruijn's strip theorem}\label{sec:debruijn}

\begin{definition}[Lean: \lean{DBN.Gen.Kernel}]
A \emph{kernel} is a continuous even function $\Phi : \R \to \R$, not identically zero, with $|\Phi(u)| \le C e^{9u}\exp(-\tfrac\pi2e^{4u})$ for $u \ge 0$. Its heat flow is $H_t(z) = \int_0^\infty e^{tu^2}\Phi(u)\cos(zu)\,du$.
\end{definition}

The template $e^{9u}\exp(-\tfrac\pi2 e^{4u})$ is the decay of Riemann's $\Phi$ in \eqref{eq:Hzeta}. The definition abstracts the \lean{dbn} island's $\Phi$ \cite{Murphy932}: the auditor instantiated the generic library with the island's $\Phi$ and recovered the island's theorem verbatim \cite[REVIEW\_GEN\_D1\_D2a \S1]{bbrecords}.

\begin{theorem}[de Bruijn \cite{deBruijn1950}; Lean: \lean{DBN.Gen.dbn\_\allowbreak{}debruijn\_\allowbreak{}parametric}]\label{thm:debruijn}
Let $\Phi$ be a kernel, and let $t_0, Y \in \R$. If every zero of $H_{t_0}$ satisfies $(\im z)^2 \le Y$, then for every $t \ge t_0$ every zero of $H_t$ satisfies $(\im z)^2 \le \max(Y - 2(t-t_0), 0)$.
\end{theorem}

\noindent
This is the form stated in \cite[Theorem 3.2]{Polymath2019} and \cite[Theorem 2.6]{FGMPS2026}, both of which attribute it to \cite[Theorem 13]{deBruijn1950}; we have not consulted de Bruijn's paper directly, and the formal proof does not depend on it. The formalized proof has four steps:

\begin{enumerate}[label=\arabic*.]
\item \emph{Discretise the heat kernel.} For $s = t - t_0 > 0$, $N \ge 1$ and $\delta = \sqrt{2s/N}$, put
\[
G_N(z) = \int_0^\infty \cosh(\delta u)^N\,e^{t_0u^2}\Phi(u)\cos(zu)\,du .
\]
Multiplying the kernel by $\cosh(\delta u)$ corresponds to the vertical shift average $(T_\delta f)(z) = \tfrac12(f(z+i\delta) + f(z-i\delta))$. So $G_N = T_\delta^N H_{t_0}$. Moreover $\cosh(u\sqrt{2s/N})^N \to e^{su^2}$, and $G_N \to H_t$ locally uniformly on $\C$ [\lean{m1W\_\allowbreak{}succ}, \lean{m1\_\allowbreak{}tendstoLocallyUniformly\_\allowbreak{}G}]. The island's design notes record that the cruder discretisation $(1 + su^2/N)^N$ preserves a zero strip but does not shrink it; this is not formalized, and that discretisation is not used.
\item \emph{The step lemma.} Let $f$ be entire, real on $\R$ and even, with an even Hadamard product. If its zeros satisfy $(\im z)^2 \le \Delta$, then the zeros of $T_\delta f$ satisfy $(\im z)^2 \le \max(\Delta - \delta^2, 0)$ [\lean{shiftAvg\_\allowbreak{}zero\_\allowbreak{}im\_\allowbreak{}sq\_\allowbreak{}le}]. The proof works factor by factor. For $\im w > 0$ beyond the threshold, each factor of the product is strictly larger at $w + i\delta$ than at $w - i\delta$. So $|f(w+i\delta)| > |f(w-i\delta)|$, which rules out $f(w+i\delta) = -f(w-i\delta)$. Iterating $N$ times lowers $\Delta$ by $N\delta^2 = 2s$.
\item \emph{Hadamard.} Each $G_N$ is even, real and entire of order at most $3/2 < 2$, so it has an even Hadamard product. The bound on the order holds for every real $t_0$, negative included, because the double-exponential decay of $\Phi$ absorbs $e^{t_0u^2}$ [\lean{m1\_\allowbreak{}evenHadamardData\_\allowbreak{}W}].
\item \emph{Hurwitz.} The $G_N$ have no zeros in the open set $\{(\im z)^2 > \max(Y - 2s, 0)\}$ and converge locally uniformly to $H_t$. By Hurwitz's theorem $H_t$ has no zeros there either, provided $H_t \not\equiv 0$ [\lean{hurwitz\_\allowbreak{}ne\_\allowbreak{}zero\_\allowbreak{}of\_\allowbreak{}entire}].
\end{enumerate}

\paragraph{Fourier injectivity in place of positivity.} For $\zeta$, the island proved $H_t \not\equiv 0$ from $\Phi > 0$, which gives $H_t(0) > 0$. A Dirichlet kernel $\Phi_\chi$ need not be positive. And $H^\chi_0(0) \ne 0$ is equivalent to $L(\tfrac12,\chi) \ne 0$, which is Chowla's conjecture and open in general. The generic library therefore uses Fourier injectivity instead [\lean{exists\_\allowbreak{}cosTransform\_\allowbreak{}ne\_\allowbreak{}zero}, \lean{H\_\allowbreak{}exists\_\allowbreak{}ne\_\allowbreak{}zero}]. Let $w$ be even, continuous, integrable and not identically zero. If its cosine transform vanished on $\R$, then, $w$ being even, its whole Fourier transform would vanish, and Fourier inversion (Mathlib's \lean{Continuous.fourierInv\_\allowbreak{}fourier\_\allowbreak{}eq}) would give $w = 0$. Applying this to $w(u) = e^{tu^2}\Phi(u)$ shows $H_t \not\equiv 0$ for every real $t$, with no arithmetic input.

\paragraph{Other decay rates.} A \emph{decay kernel} [\lean{DBN.Gen.DecayKernel}] is a continuous even $\Phi \not\equiv 0$ with $|\Phi(u)| \le Ce^{au}\exp(-be^{4u})$ for $u \ge 0$, for some real $a$ and some $b > 0$. Its rescaling $\Phi(2\,\cdot)$ is a kernel in the sense above: with $x = e^{4u}$,
\[
2au - bx^2 \le \frac{(|2a-9|+\pi/2)^2}{4b} + 9u - \frac\pi2 x \qquad (u \ge 0),
\]
and the substitution $v = 2u$ gives $H^{\Phi}_t(z) = 2H^{\Phi(2\cdot)}_{4t}(2z)$ [\lean{DecayKernel.template}, \lean{DecayKernel.Hd\_\allowbreak{}eq}]. Under the rescaling a strip $(\im z)^2 \le Y$ at time $0$ becomes $(\im z)^2 \le 4Y$, and time $t$ becomes $4t$. \Cref{thm:debruijn} then gives:

\begin{corollary}[Lean: \lean{DecayKernel.real\_\allowbreak{}zeros\_\allowbreak{}of\_\allowbreak{}strip}]\label{cor:strip}
Let $Y \ge 0$. If every zero of $H^\Phi_0$ satisfies $(\im z)^2 \le Y$, then every zero of $H^\Phi_t$ is real for $t \ge Y/2$.
\end{corollary}

\noindent
By \eqref{eq:decay}, $\Phi_\chi$ is a decay kernel with $a = 1 + 2k$ and $b = \pi/(2N)$, once it is known to be even [\lean{thetaDecayKernel}].

\subsection{From the half-plane $7/8$ to the strip $D = 3/4$}\label{sec:strip}

\begin{proposition}[Lean: \lean{DBN.Gen.zero\_\allowbreak{}strip}]\label{prop:strip}
Let $\chi$ be real, primitive and non-trivial, with $\Phi_\chi$ even, and suppose $L(s,\chi) \ne 0$ for $\re s > 7/8$. Then every zero $z$ of $H^\chi_0$ satisfies $(\im z)^2 \le 9/16$.
\end{proposition}

\begin{proof}
Let $H^\chi_0(z) = 0$ and $s = \tfrac12 + \tfrac{iz}2$, so $\re s = \tfrac12 - \tfrac12\im z$. By the bridge and $N^{(s+k)/2} \ne 0$, $\Lc_{\mathrm{M}}(s,\chi) = 0$. If $\re s > 7/8$, then $\GR(s+k) \ne 0$ and $L(s,\chi) = \Lc_{\mathrm{M}}(s,\chi)/\GR(s+k) = 0$, a contradiction; so $\re s \le 7/8$. A real-valued character is its own inverse: $|\chi(a)| = 1$ on units together with $\chi(a) \in \R$ forces $\chi(a) = \pm1$ [\lean{inv\_\allowbreak{}eq\_\allowbreak{}self\_\allowbreak{}of\_\allowbreak{}real}]. So \eqref{eq:FE} gives $\Lc_{\mathrm{M}}(1-s,\chi) = N^{s-1/2}\varepsilon(\chi)\Lc_{\mathrm{M}}(s,\chi) = 0$, and the same argument applied to $1-s$ gives $\re(1-s) \le 7/8$, that is, $\re s \ge 1/8$. Hence $|\im z| \le 3/4$.
\end{proof}

\noindent
Primitivity enters in three places: here, through \eqref{eq:FE}; in the theta functional equation \eqref{eq:thetaFE}; and in Gauss's theorem. With $Y = 9/16$, Corollary~\ref{cor:strip} gives real zeros for $t \ge 9/32 = 2(7/8 - 1/2)^2$ [\lean{DBN.Gen.dirichlet\_\allowbreak{}real\_\allowbreak{}zeros}]. Discharging evenness by \cref{thm:gauss} [\lean{DBN.Gen.real\_\allowbreak{}primitive\_\allowbreak{}real\_\allowbreak{}zeros}], and the half-plane by OpenAI's theorem with $\chi \ne 1$ excluding the pole, proves \cref{thm:main}(a).

\subsection{An example: $\chi_{-4}$}\label{sec:chi4}

The non-trivial character modulo $4$, $\chi_{-4}(n) = (-1)^{(n-1)/2}$ for odd $n$, is odd ($k = 1$), so
\[
\theta_{\chi_{-4}}(x) = \sum_{n\ \mathrm{odd}}(-1)^{(n-1)/2}\,n\,e^{-\pi n^2x/4},\qquad
\Phi_{\chi_{-4}}(u) = e^{3u}\,\theta_{\chi_{-4}}(e^{4u}).
\]
Its Gauss sum is $\tau(\chi_{-4}) = e^{2\pi i/4} - e^{2\pi i\cdot3/4} = 2i$, so $\varepsilon(\chi_{-4}) = 2i/(i\cdot2) = 1$, and $\theta_{\chi_{-4}}(1/x) = x^{3/2}\theta_{\chi_{-4}}(x)$. In Lean this character is \lstinline[basicstyle=\ttfamily,backgroundcolor={}]|chiM4 := ZMod.χ₄.ringHomComp (Int.castRingHom ℂ)|. Its primitivity, non-triviality, reality, oddness, $\tau = 2i$ and $\varepsilon = 1$ are proved by direct computation in \lean{GenChi4.lean}, independently of the general Gauss theorem [\lean{chiM4\_\allowbreak{}isPrimitive}, \lean{gaussSum\_\allowbreak{}chiM4}, \lean{rootNumber\_\allowbreak{}chiM4}]. The corresponding special case was the first one to be proved in OpenAI's environment:
\begin{lstlisting}
theorem bb_chiM4_real_zeros_oai :
    ∀ t : ℝ, 9 / 32 ≤ t → ∀ z : ℂ, DBN.Gen.Hchi DBN.Gen.chiM4 t z = 0 → z.im = 0
\end{lstlisting}
As an independent numerical sanity check of the bridge, the auditor compared both sides in \lean{mpmath} at 30 digits \cite[REVIEW\_D2b\_BRIDGE]{bbrecords}. The left side was $H^{\chi_{-4}}_0(z)$, computed by quadrature; the right side was $\tfrac18 4^{(s+1)/2}\GR(s+1)L(s,\chi_{-4})$:
\begin{center}\small
\begin{tabular}{@{}lll@{}}
\toprule
$z$ & both sides & $|\text{difference}|$ \\
\midrule
$0$ & $0.122589201757214188$ & $5\cdot10^{-31}$ \\
$1$ & $0.120217919764671738$ & $1\cdot10^{-32}$ \\
$0.5 + 0.25i$ & $0.122140362295467626 - 0.000595884686946844\,i$ & $5\cdot10^{-31}$ \\
\bottomrule
\end{tabular}
\end{center}
At $z = 0$ the value is $\tfrac18\Lc(\tfrac12,\chi_{-4})$, numerically non-zero; the proof does not use this. FGMPS prove the stronger unconditional bound $\Lc_{\chi_{-4}} \le 0.26$ \cite[\S7.5]{FGMPS2026}. Moreover, once OpenAI's theorem is granted, $9/32$ for $\chi_{-4}$ already follows on paper from FGMPS's Remark~3.4. The example is here to make the objects concrete; what is new for $\chi_{-4}$, as for every conductor, is only that the bound is a kernel-checked theorem.

\subsection{What the proof does not use}

Apart from OpenAI's theorem, whose internals we do not describe, the proof uses no numerical verification of GRH, no zero computation, no Riemann--Siegel formula, no interval arithmetic, no positivity of $\Phi_\chi$, and no non-vanishing at the central point. Its only arithmetic inputs are OpenAI's half-plane and Gauss's sign theorem. Everything else is the functional equation and complex analysis.

%======================================================================
\section{The formalization}\label{sec:formal}
%======================================================================

\subsection{Architecture}

The proof has five layers. Each imports only Mathlib and the layers above it; the composition layer also imports OpenAI's \path{OAI.NumberTheory.DirichletL.Nonvanishing}. Line counts were taken from the audited sources; the md5 prefixes are those recorded in the auditor's reviews \cite{bbrecords} and were rechecked for this draft.

\begin{center}\small
\begin{tabular}{@{}P{0.2\linewidth}P{0.11\linewidth}P{0.62\linewidth}@{}}
\toprule
Layer & Files, lines & Content \\
\midrule
\lean{dbn} island \cite{Murphy932} (commit \lean{afd94c7}) & 7, 2{,}193 & \lean{DBNStep} (the shift average and the step lemma), \lean{DBNHurwitz}, and \lean{DBNHadamard} with its four sub-modules (Hadamard factorisation of even entire functions of order $<2$). Pure complex analysis, with no $\Phi$. \\
\lean{Gen} & 13, 2{,}788 & \lean{GenFourier}, \lean{GenDefs}, \lean{GenHeatApprox}, \lean{GenM1Approx}, \lean{GenM1Parametric}: \cref{thm:debruijn} for abstract kernels. \lean{GenScale}: Corollary~\ref{cor:strip}. \lean{GenTheta}, \lean{GenMellin}, \lean{GenMellin2}, \lean{GenBridge}: \cref{sec:theta}. \lean{GenStrip}: Proposition~\ref{prop:strip}. \lean{GenThetaFE}: \eqref{eq:thetaFE}, \lean{Hchi}. \lean{GenChi4}: \cref{sec:chi4}. \\
\lean{Gauss} & 6, 1{,}419 & \lean{GaussQuad} (S1), \lean{GaussPrime} (S2), \lean{GaussCRT}, \lean{GaussPow}, \lean{GaussTwo}, \lean{GaussSign} (S3): \cref{thm:gauss} and \lean{real\_\allowbreak{}primitive\_\allowbreak{}real\_\allowbreak{}zeros}, the main theorem with the half-plane as a hypothesis. \\
\lean{GaussBridge} & 1, 36 & \lean{Hchi\_\allowbreak{}bridge}: part (b) with no extra hypotheses. \\
Composition & 2 & \lean{BbOaiRealPrim}, \lean{BbOaiMeaning}: OpenAI's theorem discharges the half-plane; the conjunction. \\
\bottomrule
\end{tabular}
\end{center}

\noindent
The seven island files (in the order of the table) have md5 prefixes \lean{e76f64ba}, \lean{1ca47174}, \lean{8069c2da}, \lean{da8f08cd}, \lean{f6e28cd1}, \lean{b4d07462}, \lean{46c3b886} (\lean{DBNHadamardCount}, \lean{DBNHadamardLinear}, \lean{DBNHadamardMean}, \lean{DBNHadamardProduct}, \lean{DBNStep}, \lean{DBNHurwitz}, \lean{DBNHadamard}). The thirteen \lean{Gen} files and six \lean{Gauss} files have md5 prefixes
\lean{261aa493}, \lean{c5c6e856}, \lean{282fe6c7}, \lean{23f22b65}, \lean{99f5dfdc}, \lean{6f608628}, \lean{3dbf6dec}, \lean{cb5c2820}, \lean{67e1af81}, \lean{5691f87d}, \lean{8aa81558}, \lean{8aeec9b0}, \lean{9acc262c} and
\lean{51b4e494}, \lean{6fb94967}, \lean{dce09dba}, \lean{16bd30ef}, \lean{3800b2cc}, \lean{03883647};
\lean{GaussBridge} is \lean{4f746054}; the composition files are \lean{BbOaiRealPrim} \lean{4c569252} and \lean{BbOaiMeaning} \lean{23ea25e9}. Four of the \lean{Gen} files (\lean{GenDefs}, \lean{GenHeatApprox}, \lean{GenM1Approx}, \lean{GenM1Parametric}) are generated mechanically from the island by a text transform that abstracts $\Phi$. The auditor re-ran the generator and obtained byte-identical outputs \cite[REVIEW\_GEN\_D1\_D2a]{bbrecords}.

The main theorem with the half-plane as a hypothesis is:
\begin{lstlisting}
theorem DBN.Gen.real_primitive_real_zeros {N : ℕ} [NeZero N] (χ : DirichletCharacter ℂ N)
    (hprim : χ.IsPrimitive) (hχ1 : χ ≠ 1) (hreal : ∀ n : ℕ, (χ n).im = 0)
    (hL : ∀ s : ℂ, (7 / 8 : ℝ) < s.re → DirichletCharacter.LFunction χ s ≠ 0) :
    ∀ t : ℝ, 9 / 32 ≤ t → ∀ z : ℂ, Hchi χ t z = 0 → z.im = 0
\end{lstlisting}
\noindent
The composition with OpenAI's theorem, in full, is:
\begin{lstlisting}
import GaussSign
import OAI.NumberTheory.DirichletL.Nonvanishing

theorem bb_real_primitive_real_zeros_oai :
    ∀ {N : ℕ} [NeZero N] (χ : DirichletCharacter ℂ N), χ.IsPrimitive → χ ≠ 1 → (∀ n : ℕ, (χ n).im = 0) →
    ∀ t : ℝ, 9 / 32 ≤ t → ∀ z : ℂ, DBN.Gen.Hchi χ t z = 0 → z.im = 0 :=
  fun χ hprim hχ1 hreal ↦ DBN.Gen.real_primitive_real_zeros χ hprim hχ1 hreal (fun _ hs ↦
    OAI.DirichletCharacter.LFunction_ne_zero_of_seven_eighths_lt_re χ hs (fun h ↦ hχ1 h.1))
\end{lstlisting}
\noindent
OpenAI's side condition $\neg(\chi = 1 \wedge s = 1)$ follows from $\chi \ne 1$.

\subsection{Two environments}

OpenAI's theorem lives at Lean \lean{v4.34.1} with Mathlib \lean{d13f23b}. Our chain was developed and first audited at Lean \lean{v4.35.0-rc3} with Mathlib \lean{0493595e}. Combining a theorem from one Mathlib revision with a hypothesis stated in another would leave a seam: the two revisions must define \lean{DirichletCharacter.\allowbreak{}LFunction} in the same way. We avoided that question. The whole chain was rebuilt from source inside OpenAI's built environment, with an empty port diff: all $27$ chain files, island included, are byte-identical to the versions audited at \lean{v4.35.0-rc3} (the island files differ from the island's release \lean{afd94c7} by two previously audited port lines). The final theorem was then proved there \cite[REVIEW\_GAUSS\_SIGN\_AND\_OAI \S B.1; ADD1 \S3]{bbrecords}. The OpenAI environment is a materialised copy of \lean{openai/\allowbreak{}math}'s \lean{lean/\allowbreak{}} directory from a replay on 7 October 2026. A recursive \lean{diff} of its \lean{OAI} tree against a git clone at \lean{adc7f1241b42}, covering 121,734 files, is empty. The lakefile, manifest, toolchain and patches agree, and \lean{Nonvanishing.lean} (md5 \lean{ea8efdb1}) re-elaborates from source \cite[REPORT\_OAI\_PORT (b)]{bbrecords}. The dependency checkouts of that local tree (\path{.lake/packages}, Mathlib among them) were not hashed against the pinned commits; the sandboxed Comparator run re-created them from the public pins: Mathlib from its official build cache at \lean{d13f23b}, and the two other packages OpenAI's theorem imports (\lean{PrimeNumberTheoremAnd}, \lean{rellich-kondrachov}) cloned at their pins and patched with OpenAI's patches. Its solution export is byte-identical to the one the auditor made from the local tree (\cref{sec:comparator}), which closes this gap for the closure of the main theorem \cite[REVIEW\_COMPARATOR\_LAMBDA\_CHI \S4]{bbrecords}.

\subsection{The trust base}\label{sec:trust}

To believe \cref{thm:main} on the strength of this formalization, one has to trust:
\begin{enumerate}
\item \emph{The Lean kernel}, plus the further checkers that replayed the export (\cref{sec:verification}), from three implementation lineages.
\item \emph{OpenAI's Lean proof} of the $7/8$ half-plane, as checked by the kernel. Its mathematical manuscript \cite{OAI7over8} is unrefereed, as is the companion $11/12$ preprint \cite{OAIqrh}. Its proof term is in the replayed environment, so the kernel checks are a check of OpenAI's proof too. But no human has refereed the argument.
\item \emph{Mathlib's definitions} at \lean{d13f23b}, in particular \lean{DirichletCharacter}, \lean{IsPrimitive}, \lean{LFunction}, \lean{completedLFunction}, \lstinline[basicstyle=\ttfamily,backgroundcolor={}]|Gammaℝ|, the complex power, \lean{tsum} and the Bochner integral, as compiled in Mathlib's official build cache at \lean{d13f23b}, from which the sandboxed run took them; the local build tree gives a byte-identical export of the main theorem's environment (\cref{sec:comparator}).
\item \emph{Our four definitions} of \cref{sec:statement} (\lean{thetaC}, \lean{PhiC}, \lean{parityK}, \lean{Hchi}). Part (b) of \cref{thm:main}, with Fourier uniqueness, is the check that \lean{Hchi} means what we say; what remains to be read is its outer heat-integral shape. The identification with FGMPS's normalisation for $t \ne 0$ (\cref{sec:normalisation}) is by hand.
\item \emph{The checking runs themselves.} The local replays on macOS were unsandboxed. The sandboxed Comparator run of \cref{sec:comparator} ran as uid 0 with all capabilities dropped; it was made by another agent session, and the auditor read all of its logs but did not repeat its kernel runs.
\end{enumerate}
Not trusted: any proof in our chain or in OpenAI's library. Every proof is checked by the kernel.

\subsection{How to read and check it}

A reader who wants to check the statement without reading proofs needs three things: \lean{ChallengeDefs.lean} (the four definitions, 56 lines, md5 \lean{67e538bc}); \lean{Challenge.lean} (the statement of \cref{sec:statement} with \lean{:= sorry}, md5 \lean{73c1bf94}); and the Mathlib definitions listed above. A reader who wants to follow the proof should read \lean{GenStrip.lean} (89 lines) first: it is the proof of Proposition~\ref{prop:strip} and of the main theorem in hypothesis form. Then read \lean{GenScale.lean} and \lean{GenM1Parametric.lean} for de Bruijn, and \lean{GenThetaFE.lean} and \lean{GaussSign.lean} for the arithmetic. The release repository's README gives two routes for reproducing the result, summarised here.
\begin{enumerate}
\item \emph{The auditor's build.} Clone \lean{openai/\allowbreak{}math} and check out \lean{adc7f124}; its Lean project is \path{lean/}, with toolchain \lean{v4.34.1} and Mathlib \lean{d13f23b}. In \path{lean/}, run \lean{lake update}, \lean{lake exe cache get} and \lean{lake build OAI.NumberTheory.DirichletL.Nonvanishing}, which is OpenAI's documented procedure and applies their dependency patches. Afterwards, check that \path{lean/lake-manifest.json} is unchanged (md5 \lean{8c3ec864}), for example with \path{<release>/scripts/check_manifest.sh}~\path{<clone>/lean/lake-manifest.json}; otherwise the dependencies are no longer those pinned by OpenAI. Then compile the $29$ chain modules of the release (\path{lean/island}, \path{lean/gen}, \path{lean/gauss}, \path{lean/main}) with plain \lean{lean}, in the order listed in the README, against that build. Finally run \lean{BbOaiMeaningCheck.lean} (axioms and closure walks for both halves), \lean{BbOaiMeaningRef.lean} (the \lean{pp.all} statement with no OpenAI module in scope) and \lean{BbOaiMeaningTamper.lean}, which must fail with $3$ errors, and compare the outputs with \path{logs/oai_env/}. This mirrors how the Lean results of \cref{sec:verification} were obtained; the auditor used OpenAI's tree as built in the 7 October replay (\cref{sec:formal}) rather than a fresh clone.
\item \emph{Through Lake.} \path{scripts/verify_md5.sh} checks every file under \path{lean/}, \path{logs/} and \path{provenance/} against its manifest, and the challenge files against the auditor's two manifests in \path{challenge/}. \path{scripts/materialize.sh} sets up \lean{openai/\allowbreak{}math} at the pin with the chain inside it as five \lean{lean\_lib} targets; one then builds those targets and runs the same three checks with \lean{lake env lean}. This route, and the continuous-integration workflow that runs it on Linux, had not been run at the time of this version. (OpenAI's package sets \lean{autoImplicit} off; the stanza sets it back on for each of the five chain libraries, matching how the chain was compiled, but this override is untested. A mismatch could only fail loudly, as an unknown-identifier error, and could not change a statement.)
\end{enumerate}
The five-checker replay of \cref{sec:verification} exports the environment with \lean{lean4export} and runs each checker on the export, as in the scripts under \path{logs/five_checker/}.

%======================================================================
\section{Verification}\label{sec:verification}
%======================================================================

All results in this section were produced by the auditor (an AI agent session; see \cref{sec:ai}) and are recorded, with md5 sums for every input and output, in its reviews \cite{bbrecords}.

\subsection{Axioms}

\lean{\#print axioms} reports \lean{[propext, Classical.choice, Quot.sound]} for \lean{bb\_\allowbreak{}real\_\allowbreak{}primitive\_\allowbreak{}meaning\_\allowbreak{}oai}, \lean{Hchi\_\allowbreak{}bridge}, \lean{bb\_\allowbreak{}real\_\allowbreak{}primitive\_\allowbreak{}real\_\allowbreak{}zeros\_\allowbreak{}oai} and \lean{rootNumber\_\allowbreak{}eq\_\allowbreak{}one\_\allowbreak{}of\_\allowbreak{}real\_\allowbreak{}primitive}. \lean{\#print axioms} reads a cached summary, so the auditor also walked the term-level closure directly: types, values, constructors and recursor rules. The walk over the $126{,}656$ constants in the closure of the main theorem finds the same three axioms \cite[ADD1 \S3]{bbrecords}. A static scan of the Gauss and bridge sources finds no \lean{sorry}, \lean{admit}, \lean{axiom}, \lean{native\_\allowbreak{}decide}, \lean{implemented\_\allowbreak{}by}, \lean{extern}, \lean{unsafe}, \lean{opaque}, \lean{set\_\allowbreak{}option}, \lean{macro}, \lean{elab} or \lean{\#eval}. The only textual hit, in a doc comment, is the word ``admit'' \cite[\S A.2; ADD1 \S1]{bbrecords}.

\subsection{Five checkers from three lineages, and a control they must reject}

The environment containing the main theorem was exported with \lean{lean4export} \cite{lean4export} (v4.34.1 build, export format 3.1.0) to a file of $1{,}730{,}561{,}283$ bytes (md5 \lean{6888b248}). The export was given to five checkers. These are the checkers bundled with the Lean \lean{v4.35.0-rc3} toolchain, which read the same export format. Alongside it, a control was exported and checked: a ``theorem'' \lean{BbNeg.bogus : 4 = 5} admitted into the environment by switching off the elaborator's kernel check (\lean{debug.skipKernelTC}; export md5 \lean{bf973608}). A checker that accepts the control is not checking anything. The runs were guarded so that an empty or failed export aborts the run, and any acceptance of the control marks it invalid. The guards exist because an earlier run in the same programme exported zero bytes because of a toolchain mismatch, and every checker then ``accepted'' the empty file, the control included \cite[REVIEW\_CTFREE\_KERNELS]{bbrecords}.

\begin{center}\small
\begin{tabular}{@{}P{0.3\linewidth}P{0.36\linewidth}P{0.26\linewidth}@{}}
\toprule
Checker & main theorem (\lean{6888b248}) & control \lean{4 = 5} \\
\midrule
\lean{leanchecker} (Lean's kernel) & accept, 270 s & reject (rc 1) \\
\lean{leanchecker-paranoid} & accept, 320 s & reject (rc 1) \\
\lean{nanoda} \cite{nanoda} & accept, 57 s & reject (rc 101, panic in \lean{def\_\allowbreak{}eq}) \\
\lean{con-leche} & accept, ``accepted 124635 declarations (\texttt{-{}-verified})'' & reject (rc 1) \\
\lean{con-ron} & accept, ``accepted 124635 declarations (\texttt{-{}-verified})'' & reject (rc 1) \\
\lean{lean4lean} \cite{lean4lean} & N/A: crashes before reaching any target & N/A \\
\bottomrule
\end{tabular}
\end{center}

\noindent
All five accept, and all five reject the control \cite[ADD1 \S4]{bbrecords}. The rc3 \lean{lean4lean} cannot read v4.34.1 exports: it fails with ``unknown constant \lean{Decidable.intro}'' on every file, the control included, so it is recorded as not applicable rather than as a rejection. The two \lean{leanchecker} variants run Lean's own kernel. \lean{nanoda} is an independent implementation in Rust. \lean{con-leche} is a separate kernel implementation written in Lean, with its own term datatypes (\lean{ConLeche.Expr}, \lean{ConLeche.Level}, \lean{ConLeche.Name}); its help text states a soundness theorem for its default \lean{-{}-verified} mode, ``\lean{ConLeche.no\_\allowbreak{}proof\_\allowbreak{}of\_\allowbreak{}False}'' (if the declaration fold accepts, no constant of type \lean{False} is in the environment). \lean{con-ron} describes itself as a Rust port of \lean{con-leche}. Both are bundled with the Lean \lean{v4.35.0-rc3} toolchain, whose \lean{lake comparator -{}-paranoid} runs them alongside \lean{leanchecker-paranoid}, \lean{lean4lean} and \lean{nanoda} (\path{Lake/CLI/Check.lean} in the toolchain's sources). \lean{con-leche} \cite{conleche} was implemented, and proved consistent in Lean, by Claude under the supervision of Joachim Breitner at the Lean FRO; its README states the theorems \lean{no\_False\_declaration} and \lean{model\_exists} (every accepted environment has a model in set theory, following Carneiro's model of Lean's type theory). \lean{con-ron} \cite{conron} is a Rust port of it, written by AI under Breitner's supervision, with an Aeneas proof that it refines \lean{con-leche}'s checking algorithm; it differs in term representation, caching, compiler and runtime. Both ship with Lean \lean{v4.35.0} \cite{deMoura2026}. They are independent of Lean's C++ kernel, but not of each other. The five acceptances therefore come from three implementation lineages: Lean's kernel, \lean{nanoda}, and the \lean{con-leche}/\lean{con-ron} pair. Separate runs of the same kind accepted part (a) alone (export md5 \lean{ae0c51bd}, $124{,}633$ declarations) and the $\chi_{-4}$ case (\lean{7c142b2c}, $124{,}395$ declarations), and rejected the control \cite[REVIEW\_GAUSS\_SIGN\_AND\_OAI \S B.6; REPORT\_OAI\_PORT (a)]{bbrecords}.

\subsection{Tampered statements}

Each headline statement was also compiled with one number changed, with the proof unchanged. We distinguish two kinds of rejection. A \emph{genuine} rejection means the changed statement is false, or contradicts a proved lemma. A \emph{proof-reach} rejection means only that the given proof does not reach the changed statement; it shows a hypothesis or constant is load-bearing, not that the changed statement is false.

\begin{center}\small
\begin{tabular}{@{}P{0.48\linewidth}P{0.24\linewidth}P{0.2\linewidth}@{}}
\toprule
Change & Result & Kind \\
\midrule
half-plane $7/8 \to 3/4$ in the composition & rejected at OpenAI's theorem & proof-reach (OpenAI's theorem does not reach $3/4$) \\
heat time $9/32 \to 1/4$ & rejected & proof-reach \\
bridge constant $1/8 \to 1/4$ & rejected & genuine \\
parity $k \to 1-k$ in the bridge & rejected & genuine \\
bridge without primitivity & rejected at \lean{assumption} & proof-reach \\
$\varepsilon(\chi) = -1$ for real primitive $\chi$ & rejected & genuine \\
$(1+i) \to (1-i)$ in the closed form of $G(N)$ & rejected & genuine \\
$\varepsilon(\chi_{-4}) = -1$; $\tau(\chi_{-4}) = -2i$ & rejected & genuine \\
\bottomrule
\end{tabular}
\end{center}

\noindent
The auditor's earlier records called the first row ``genuine'', meaning that the rejection happens at OpenAI's interface; under the definition above it is proof-reach, since the changed hypothesis is implied by GRH and is not known to be false. Sources: \cite[REVIEW\_GAUSS\_SIGN\_AND\_OAI \S\S A.6, B.5; ADD1 \S\S2, 3; REVIEW\_D2bp\_CHI4]{bbrecords}.

\subsection{Statement identity with and without OpenAI's code}

Loading OpenAI's library into the environment could, in principle, change what a statement means, for instance through an instance or a declaration that shadows a Mathlib name. To rule this out, the auditor printed the statements with \lean{pp.all}, in which every implicit argument and instance is explicit, in two environments: one containing OpenAI's library, and one containing only our chain and Mathlib. The text of part (a) ($4{,}017$ characters) and the text of part (b) ($7{,}778$ characters) are byte-identical across the two \cite[REVIEW\_GAUSS\_SIGN\_AND\_OAI \S B.3; ADD1 \S3]{bbrecords}. At the $\chi_{-4}$ stage the OpenAI environment had $13{,}708$ modules, with no module-name collisions and no \lean{DBN.*} constant from a foreign module \cite[REPORT\_OAI\_PORT]{bbrecords}.

\subsection{Independent numerics}

These computations check conventions and normalisations; they are plain floating point or \lean{mpmath}, not certificates. With its own implementation of the Kronecker symbol, in Mathlib's normalisation $\varepsilon = \tau/(i^k\sqrt N)$, the auditor checked \cite[\S A.7]{bbrecords}:
\begin{itemize}
\item the closed form of $G(N)$ for $N = 1, \dots, 200$ (largest error $6.6\cdot10^{-38}$);
\item $\varepsilon = 1$ for real characters of odd prime conductor $p < 500$ ($4.2\cdot10^{-41}$);
\item $\varepsilon(\chi_D) = 1$ for every fundamental discriminant $0 < |D| \le 2000$, of both signs: $1{,}218$ discriminants, largest error $4.3\cdot10^{-31}$.
\end{itemize}
The bridge was also checked numerically for $\chi_{-4}$ (\cref{sec:chi4}).

\subsection{The sandboxed Comparator run}\label{sec:comparator}

\noindent
\emph{What was run} \cite[COMPARATOR\_RUN\_SPEC]{bbrecords}. Comparator \cite{Comparator} checks a solution against a separately written challenge. It elaborates the challenge statement against a minimal trusted definitions file. It requires that the solution's theorem have exactly that type and that every definition the statement uses be identical, as a \lean{ConstantInfo}, to the challenge's. It then checks the axiom whitelist and replays the proof through the kernel and further checkers, inside a sandbox. For this theorem the auditor wrote:
\begin{itemize}
\item \lean{Challenge.lean} (md5 \lean{73c1bf94}): the statement of \cref{sec:statement}, verbatim, with \lean{:= sorry};
\item \lean{ChallengeDefs.lean} (md5 \lean{67e538bc}): the trusted definitions. It introduces exactly six non-Mathlib constants: \lean{thetaC}, \lean{PhiC} with its two auxiliary proof terms, \lean{parityK} and \lean{Hchi};
\item \lean{config.json} (md5 \lean{d563faf7}): axioms \lean{propext}, \lean{Quot.sound}, \lean{Classical.choice}; nanoda enabled;
\item two tampered challenges that must be rejected: heat time $9/32 \to 1/4$ (\lean{d6ef5198}) and bridge constant $1/8 \to 1/4$ (\lean{31163137}).
\end{itemize}
The run was made inside a \lean{bwrap} sandbox on a Linux machine, on OpenAI's environment plus our chain at the audited md5s, once with Comparator's default kernels and once with the paranoid kernel set; the auditor then read the run's logs in full and verified them \cite[REVIEW\_COMPARATOR\_LAMBDA\_CHI and its Addendum~1]{bbrecords}.

\medskip\noindent
\emph{Caveat that the result carries.} The run was made as uid 0 with all capabilities dropped (\lean{-{}-cap-drop ALL}), in an ephemeral single-tenant container; Comparator's assumption that the checker is not run by a privileged user is met in spirit, not in letter. Mathlib was taken from its official build cache at \lean{d13f23b}. The challenge-side oleans were compared with a fresh rebuild after the run, because the queued runs share one build directory (Comparator's assumption that the solution has not been compiled before the challenge).

\medskip\noindent
\emph{Results.} The run was made on 30 October 2026 by another agent session, with Comparator \lean{v4.34.0} and a \lean{bwrap} sandbox in place of \lean{landrun} (uid 0, all capabilities dropped, no network), on OpenAI's environment rebuilt from its public pins (Mathlib from its official build cache at \lean{d13f23b}; \lean{PrimeNumberTheoremAnd} and \lean{rellich-kondrachov} cloned at their pins and patched with OpenAI's patches; OpenAI's sources unmodified) and our chain at the audited md5s. With the challenge files listed above, Comparator printed ``Your solution is okay!'' both with \lean{nanoda} and Lean's kernel (log md5 \lean{fc80db0a}) and with the paranoid set \lean{leanchecker-paranoid}, \lean{nanoda}, \lean{con-leche}, \lean{con-ron} and Lean's kernel (log md5 \lean{fb164675}; \lean{con-leche} and \lean{con-ron} report $124{,}639$ declarations); \lean{lean4lean} cannot read \lean{v4.34.1} exports and was not run. In Comparator \lean{v4.34.0} that message is printed only after the statement comparison, the axiom check and every kernel have passed. Both tampered challenges were rejected with ``Challenge and solution theorem statement do not match'' (log md5s \lean{78a706f6}, \lean{d1f16445}). The auditor read all of the run's logs (\cite[REVIEW\_COMPARATOR\_LAMBDA\_CHI, Addendum~1]{bbrecords}) and reproduced the solution export independently, from its own build: it is byte-identical ($1{,}730{,}282{,}328$ bytes, md5 \lean{5ca44ad5}), so the sandboxed kernels checked, byte for byte, the export of the auditor's own build. The five-checker replay of \cref{sec:verification} used an export of that build with fewer targets (md5 \lean{6888b248}): its $124{,}635$ declaration records contain the whole closure of the main theorem and are a subset of these $124{,}639$. The extra four, \lean{Nat.lor}, \lean{Nat.xor}, \lean{String.mk} and an auxiliary proof, lie outside that closure; Comparator exports them because it lists them as kernel primitives. The auditor also ran \lean{con-leche} on its own copy of the \lean{5ca44ad5} export, and it accepted all $124{,}639$ declarations \cite[REVIEW\_COMPARATOR\_LAMBDA\_CHI, Addendum~2]{bbrecords}. The solution was compiled inside the sandbox, in an earlier attempt of the same Comparator command that a container restart interrupted (log md5 \lean{144a7c6b}); the recorded run replayed those build outputs. Before that attempt the build directory was wiped and Mathlib's cache downloaded afresh, discarding an earlier attempt that had been made without the capability drop. Between the two attempts one unsandboxed \lean{lake build Challenge} was run to test Lake's replay behaviour; it compiled nothing (these preparatory steps are from the runner's record; the logs are consistent with them). Because the queued runs share one build directory, the challenge-side oleans were rebuilt afterwards in a fresh directory and found identical (Comparator's assumption that the solution was not compiled before the challenge) \cite[REVIEW\_COMPARATOR\_LAMBDA\_CHI \S\S2--5 and Addenda~1--2]{bbrecords}.

%======================================================================
\section{Optional: companion results for $\zeta$ that use no OpenAI code}\label{sec:companion}
%======================================================================

\emph{This section is independent of the rest of the paper and may be omitted.}

Cureton and Thorner \cite{CuretonThorner2026} proved, with the AxiomProver system, a zero-free half-plane for every Dedekind zeta function. Version~1 of their report states $\re s > 1 - 1/6000$, as does their repository \cite{CuretonThornerLean} at its first public commit \lean{284348a4}. Commit \lean{201bc40} (13 October 2026) strengthens the theorem \lean{main} to $\re s > 31/32$ for every number field; we use the repository at \lean{09afbb4}, which also adds a Hecke-$L$ corollary at $31/32$. We have seen no revised report stating $31/32$; that half-plane rests on the Lean repository. At $K = \mathbb{Q}$ this gives $\zeta(s) \ne 0$ for $\re s > 31/32$, $s \ne 1$. The reduction from their statement to Mathlib's \lean{riemannZeta} was written and audited in this programme \cite[REVIEW\_CTQ]{bbrecords}. Their code imports no module of \lean{openai/\allowbreak{}math} \cite[REVIEW\_CTQ]{bbrecords}. Two classical consequences were then formalized:

\begin{enumerate}
\item \emph{$\Lc \le 225/512$.} For the $\zeta$ heat flow \eqref{eq:Hzeta}, as defined on the \lean{dbn} island: every zero of $H_t$ is real for $t \ge 225/512 = 2(31/32 - 1/2)^2$. This is the same composition as Proposition~\ref{prop:strip}, with $D = 15/16$.
\item \emph{Ingham's theorem in the kernel.} The general statement \lean{Ingham.ingham}:
\begin{lstlisting}
∀ {a : ℝ}, 1/2 ≤ a → a < 1 → (∀ s : ℂ, a < s.re → riemannZeta s ≠ 0) →
  ∀ ε : ℝ, 0 < ε → ∃ C, ∀ x : ℝ, 2 ≤ x → |ψ x − x| ≤ C * x ^ (a + ε)
\end{lstlisting}
with Mathlib's Chebyshev $\psi$ and real powers, and $C$ depending only on $a$ and $\varepsilon$. The proof uses a Perron formula with the weight $1/(s(s+1))$. The main term comes from writing $-\zeta'/\zeta = 1/(s-1) - \zeta_1'/\zeta_1$ with $\zeta_1 = (s-1)\zeta$ and evaluating the $1/(s-1)$ piece exactly. The rest of the integrand, $(\zeta_1'/\zeta_1)(s)\,x^s/(s(s+1))$, is holomorphic on $\re s > a$, $s = 1$ included, so the shift from $\re s = 2$ to a line $\re s = \sigma_0$ with $a < \sigma_0 < 1$ needs no residue theorem. A Borel--Carath\'eodory bound controls $\zeta_1'/\zeta_1$. At $a = 31/32$ it gives $\psi(x) = x + O_\varepsilon(x^{31/32+\varepsilon})$. This is the classical direction of Ingham \cite{Ingham1932}; nothing in it is new mathematics.
\end{enumerate}

\noindent
Neither bound is close to the state of the art: $225/512 \approx 0.439$ is far weaker than Platt and Trudgian's $\Lc \le 0.2$ \cite{PlattTrudgian2021} and Polymath15's $0.22$ \cite{Polymath2019}, and $\Lc \le 9/32$ for $\zeta$ has been kernel-checked from OpenAI's theorem \cite{Murphy932}. Their only point is independence: both hold in an environment that contains no OpenAI code.

\emph{Status of the checks.} In hypothesis form, with Cureton--Thorner's conclusion $\forall s,\ 31/32 < \re s \to s \ne 1 \to \zeta(s) \ne 0$ as an explicit hypothesis, both compositions carry the standard three axioms. Six checkers (\lean{leanchecker}, \lean{leanchecker-paranoid}, \lean{lean4lean}, \lean{nanoda}, \lean{con-leche}, \lean{con-ron}) accepted both, and all six rejected the same kind of \lean{skipKernelTC} control. This is a same-version run, in which \lean{lean4lean} can read the export \cite[REVIEW\_CTFREE\_KERNELS]{bbrecords}. Tampers (the exponent $15/16$ from the $31/32$ half-plane, and the quantifier swap $\exists C\,\forall\varepsilon$) were rejected \cite[REVIEW\_INGHAM]{bbrecords}. The unconditional forms rest on Cureton--Thorner's theorem \lean{main}. The auditor did not run their code. All three Comparator runs below were made on 30--31 October 2026 by another agent session. Each used the toolchain-provided \lean{lake comparator} of Lean \lean{v4.35.0-rc3} with \lean{-{}-paranoid} (\lean{leanchecker-paranoid}, \lean{lean4lean}, \lean{nanoda}, \lean{con-leche}, \lean{con-ron} and Lean's kernel), in a \lean{bwrap} sandbox as uid 0 with all capabilities dropped, on Cureton--Thorner's repository at \lean{09afbb4} (its lakefile extended by 30 additive lines that only declare the libraries of our compositions and challenges) with Mathlib from its official build cache. The installed toolchain, checkers included, is file-for-file the published \lean{v4.35.0-rc3} release. The auditor checked each run's relayed configuration, challenge and solution md5s, sandbox probe, verdicts and exit code against its own records, but did not see the logs \cite[REVIEW\_CT\_RERUNS]{bbrecords}. Cureton--Thorner's own challenge (\lean{main} and \lean{heckeCorollary}): ``Your solution is okay!'', all six kernels accepting (log md5 \lean{a33d06b0}). The P2 and Ingham statements are kernel-checked from the type of \lean{main} itself, whose reduction to Mathlib's \lean{riemannZeta} was audited separately \cite[REVIEW\_CTQ, Addendum~1]{bbrecords}. P2: PASS, six kernels (log \lean{fed70f9a}); tamper ($224/512$) rejected, exit 1 (log \lean{90424a05}); the trusted definitions of the P2 challenge match the audited \lean{DBNDefs} (md5 \lean{acfa852d}). Ingham: PASS, six kernels (log \lean{dd6fa95b}); tamper ($15/16$) rejected (log \lean{33d0ea1c}). Every solution module these runs used was compiled inside the sandbox, except Cureton--Thorner's vendored class-field-theory package, which had been compiled earlier without it; rebuilt inside the sandbox at the same path, it gave byte-identical build outputs. These runs replace earlier ones whose sandbox did not drop capabilities. The challenge-side build outputs were rebuilt afterwards in a fresh directory and found byte-identical (Comparator's assumption that the solution was not compiled before the challenge). The result rests on Cureton--Thorner \lean{main} at \lean{09afbb4}.

%======================================================================
\section{Version 2: an unconditional large-conductor bound, and every conductor below $9/32$}\label{sec:v2}
%======================================================================

This section is new in version~2. It records three bounds on $\Lc_\chi$ obtained between 31 October and 1 November 2026 by AI agent sessions in the same programme (\cref{sec:ai}), after version~1 was published, and, from version 2.1.0, a fourth (\cref{sec:W2B3}) obtained on 2--3 November 2026 by the same kind of session. \emph{Nothing in it is kernel-checked unless it says so.} The records cited as \cite{dq1records} are listed in Appendix~\ref{app:v2records} and archived with their md5s in \path{provenance/}; the certificates are in \path{certificates/} of the release. We have tried to make the section readable without the records, but every number in it comes from them, and a reader who wants to check a number should go to the ledger named beside it.

\paragraph{Conventions for this section.} The records write the conductor as $m$ and $L$ for $\log(m/\pi)$. Here we keep $N$ for the conductor and write $\ell = \log(N/\pi)$, to avoid a clash with $L(s,\chi)$; in Lean the same quantity is \lean{C2.Lc m}. The parity is $a = k(\chi) \in \{0,1\}$. The heat flow is FGMPS's $H^L_t$, in its Gaussian form \cite[(4.1)]{FGMPS2026}, read at source:
\begin{equation}\label{eq:gaussform}
H_t(z) = \int_{-\infty}^{\infty} \xi_L\Bigl(\frac{1+iz}{2} + \sqrt t\,v\Bigr)\frac{e^{-v^2}}{\sqrt\pi}\,dv ,
\end{equation}
where $\xi_L(s) = (N/\pi)^{(s+a)/2}\Gamma(\tfrac{s+a}{2})L(s,\chi)$ is the completed $L$-function, our $\Lc(s,\chi)$. By \cref{sec:normalisation}, $H^L_t = 8H^\chi_t$ for every real $t$, so the zeros, and the constant $\Lc_\chi$, are the same as in \cref{thm:main}; that identification is by hand, not formal. All statements below are of the form ``every zero of $H_t$ is real for every $t \ge c$'', as in Remark~\ref{rem:sinf}.

\subsection{The method: a floor for the logarithmic derivative, and one contraction step}\label{sec:v2method}

\paragraph{The quantity.} For $x \in \R$, $p > 0$ and $t$ real, put
\[
S_\chi(x,p,t) \;=\; -\im\,\frac{H'_t}{H_t}(x+ip) .
\]
When the zeros of $H_t$ are all below height $p$, this is a sum over the zeros of positive terms, and it measures how strongly the point $x + ip$ ``feels'' the zeros beneath it.

\paragraph{The contraction inequalities (cited; stated for $\xi$).} Planat \cite[Theorem 2.1]{Planat2026} proves, for Riemann's $\xi$ in the Polymath normalisation, that along every interval of time on which the maximal-height zero $z(t) = x(t) + iy(t)$ is simple,
\[
\frac{d}{dt}\,y(t)^2 \;\le\; -2 - 4y(t)^2\,G_t(z(t)),
\]
where $G_t(z) \ge 0$ collects the interaction of $z$ with the other zeros; dropping $G$ gives de Bruijn's classical $(y^2)' \le -2$. His Proposition~2.3 then bounds $G$ from below by the logarithmic derivative at a point above the zero: for $\eta \ge \sqrt5\,y$, the point $x + i\eta$ is zero-free and
\[
G_t(z) \;\ge\; \frac{L_\eta(t;x)}{\eta} - \frac{2}{\eta^2 - y^2},\qquad L_\eta(t;x) = -\im\frac{H'_t}{H_t}(x+i\eta).
\]
Hakimoglu-Brown \cite[p.~9]{HakimogluBrown2026} writes the combination, for $q = y^2$ and $p^2 \ge 5q$, as
\begin{equation}\label{eq:HBineq}
q' \;\le\; -2 - \frac{4q}{p}\,S(x,p,t) + \frac{8q}{p^2 - q},
\end{equation}
and his Theorem~9 turns it into a stepped comparison: if $S(x,p,t) \ge C$ on a time interval, for all relevant $x$, and $p^2 \ge 5q$ there, then with $k = 4C/p - 8/(p^2 - q) > 0$ one has $q' \le -2 - kq$, so the squared height of the highest zero lies under the barrier $Q(t) = (q_0 + 2/k)e^{-k(t - t_0)} - 2/k$ once it starts at $q_0$ at time $t_0$. Both papers are about $\zeta$: Planat's uses the Polymath normalisation and Hakimoglu-Brown's Theorem~9 carries constants specific to his application ($t_n \le 1/5$, a floor $b = 1/25$, and a finite region beyond which his Lemma~2 confines the zeros). What is used here is a single step of that comparison, transposed to $H_t$ for a Dirichlet character with its own confinement lemma and its own constants. The transposition, with the treatment of multiple zeros and of zeros on $i\R$ through FGMPS's Proposition~3.2 and Lemma~3.7, is an AI-written argument in class \ref{T3} \cite[P5 memo \S A.13.3; Y3 memo \S3.2]{dq1records}, reviewed by the auditor.

\paragraph{The landing time.} Suppose the zeros of $H_0$ lie in $|\im z| \le \sqrt{q_0}$. On $[0, t_{\min}]$ the classical inequality gives $q(t_{\min}) \le q_0 - 2t_{\min}$. If $S_\chi(x,p,t) \ge C$ for all $x$ and all $t \in [t_{\min}, T]$, with $p^2 \ge 5q_0$ and $k = 4C/p - 8/(p^2 - q_0) > 0$, then the barrier reaches a small level $b > 0$ at time
\begin{equation}\label{eq:landing}
T \;=\; t_{\min} + \frac1k\log\frac{2 + k(q_0 - 2t_{\min})}{2 + kb},
\end{equation}
and de Bruijn's strip theorem, in the time-shifted form of \cite[Theorem 2.6]{FGMPS2026} ($\Lc_L \le t_0 + \tfrac12 y_0^2$), finishes: $\Lc_\chi \le T + b/2$. The comparison needs the first contact of the barrier with a zero to be attained, which is a compactness statement: a \emph{confinement lemma} saying that for $|x|$ beyond an explicit $X$ the zeros of $H_t$ at heights between $\sqrt b$ and $1$ are absent. That lemma (``Lemma Z$'$'' in the records) is derived from FGMPS's effective approximation \cite[Theorem 4.3, Proposition 4.18]{FGMPS2026}, with its constants certified in Arb; see \cref{sec:C2}.

Two features of \eqref{eq:landing} explain everything that follows. From the trivial strip $q_0 = 1$, beating $9/32$ needs $k \ge 3.765$ (a float computation, \cite[Y3 memo \S2]{dq1records}), which forces $\ell$ large: this is why the unconditional \cref{thm:C2} starts at $N = 10^8$ (and why going lower, in \cref{sec:W2B3}, takes many steps and a sharper kernel), and the obstruction is Siegel-shaped, since a Siegel zero sits on $i\R$ at height at most $1$ and is exactly what keeps $q_0 = 1$ unconditional. From OpenAI's strip $q_0 = 9/16$, on the other hand, \emph{any} $k > 0$ gives $T + b/2 < q_0/2 = 9/32$: a strict improvement on version~1 needs only a positive floor, and that is \cref{thm:Y3}.

\paragraph{The floor, and why it is uniform in $\chi$.} Write $z = x + ip$ in \eqref{eq:gaussform} and use the functional equation $\xi_L(s) = \xi_L(1-s)$ (for real $\chi$, $\varepsilon(\chi) = 1$ by \cref{thm:gauss}) to move the argument to the right of the critical line; the conductor factor $(N/\pi)^{\sqrt t v/2} = e^{\ell\sqrt t v/2}$ then completes a square in the Gaussian and shifts its centre. The result (\cite[P5 memo \S A.3, \S A.6]{dq1records}, reviewed in REVIEW\_P5) is the exact identity
\begin{equation}\label{eq:tilt}
S_\chi(x,p,t) \;=\; \frac{\ell}{4} + \frac12\,\re\frac{D'(c)}{D(c)},\qquad
D(c) = \mathbb{E}_g\bigl[\Phi_a(c + \tau g)\bigr],\quad \Phi_a(w) = \Gamma\Bigl(\frac{w+a}2\Bigr)L(w,\chi),
\end{equation}
with $g$ a standard Gaussian, $\tau = \sqrt{t/2}$, $c = R - ix/2$ and $R = \tfrac{1+p}{2} + \tfrac{t\ell}{4}$. So the conductor enters $S_\chi$ explicitly only through the term $\ell/4$, and otherwise only through the centre $R$ of the Gaussian average, which $\ell$ moves to the right; what remains is the logarithmic derivative of a Gaussian average of $\Gamma\cdot L$ along a horizontal segment centred at real part $R \ge \tfrac{1+p}{2} > 1$. In that region $L(w,\chi)$ is controlled by its Euler product alone: $\zeta(2\sigma)/\zeta(\sigma) \le |L(\sigma+it,\chi)| \le \zeta(\sigma)$ and $|L'/L(w,\chi)| \le -\zeta'/\zeta(\re w)$ \cite[Lemma 2.1 and the proof of Proposition 2.2]{BMOR2020}, with Rademacher's convexity bound \cite[Lemma 5.6, (5.3)]{BMOR2020} and the functional equation for the Gaussian tails that reach into the strip; the $\Gamma$ factor is handled by the digamma series \cite[5.7.6, 5.15.1]{DLMF} and $|\Gamma(x+iy)| \le |\Gamma(x)|$ \cite[5.6.6]{DLMF}, and $\zeta$, $-\zeta'$ and $-\zeta'/\zeta$ are decreasing on $(1,\infty)$ \cite[p.~11]{HakimogluBrown2026}. None of these inputs distinguishes one real character from another. The certified floor is therefore a statement about every real primitive $\chi$ with a given $(\ell, a)$ at once, and the zeros of $L(s,\chi)$, Siegel zeros included, never enter it. This is the sense in which the method is blind to GRH: it bounds $\Lc_\chi$ without knowing anything about where the zeros are, and that is also why it cannot say anything about $\Lc_\chi \le 0$.

\paragraph{The certifiers.} The floor is certified over boxes in the parameters $(R, \tau, x)$. Three programs cover the domain: a box certifier for $0 \le x \le 10^4$ (``G1''), which bisects boxes until each one certifies $S_\chi - \ell/4 \ge$ the target; a certifier for all $x \ge 10^4$ at once (``LX''), which evaluates the bound at $x = 10^4$ and relies on a proved monotonicity in $x$ of every piece of it (\cite[P5 memo \S A.9.2]{dq1records}; the auditor spot-checked part (a) of that lemma and read parts (b)--(g) without fully re-deriving them); and an analytic lemma for the small-$\tau$ region, which is where every $\ell > 4\cdot10^4$ lives (``STL'', \cite[P5 memo \S A.9.3]{dq1records}, with constants in Arb). Negative $x$ is reduced to positive $x$ by $S_\chi(-x,p,t) = S_\chi(x,p,t)$, which follows from $H_t(-\bar z) = \overline{H_t(z)}$ for real $\chi$; this reflection is one of the statements proved in Lean (\cref{sec:v2lean}), under an explicit reflection hypothesis on $L(s,\chi)$. Every certifier is a \lean{python-flint} program working in balls with outward rounding, run with the BLAS pinned to one thread per worker (one core per worker; the G1 cover used six workers), and every run writes a ledger with the box count, the worst certified value and the list of failures; a run with a non-empty failure list certifies nothing.

\subsection{Theorem C2: every conductor from $10^8$ on, unconditionally}\label{sec:C2}

\begin{theorem}[C2; classes \ref{T2}--\ref{T4}, not kernel-checked except where stated]\label{thm:C2}
Let $\chi$ be a real primitive Dirichlet character of conductor $N \ge 10^8$, of either parity, and let $\ell = \log(N/\pi)$. Put
\begin{equation}\label{eq:C2bound}
\begin{gathered}
t_{\min} = \min(10^{-4}, 4/\ell),\qquad b = \min(10^{-4}, \ell^{-2}),\qquad k = \frac{4}{\sqrt5}\Bigl(\frac{\ell}{4} - 1\Bigr) - 2,\\[2pt]
T(\ell) = t_{\min} + \frac1k\log\frac{2 + k(1 - 2t_{\min})}{2 + kb}.
\end{gathered}
\end{equation}
Then every zero of $H_t$ is real for every $t \ge T(\ell) + b/2$. That is, $\Lc_\chi \le T(\ell) + b/2$.
\end{theorem}

\noindent
\emph{Hypotheses.} Only that $\chi$ is real and primitive with $N \ge 10^8$. GRH is not assumed; no zero-free half-plane is assumed (the initial strip is the trivial $|\im z| \le 1$, so OpenAI's theorem is not used); no numerical verification of zeros is used. A Siegel zero, if there is one, lies on $i\R$ at height at most $1$ and is untouched by the argument.

\emph{Values} (Arb, \cite[P5 memo \S A.13.1]{dq1records}; REVIEW\_P5\_A12 item 7). At $N = 10^8$ the bound is $0.2764281989$ (the ball has radius $10^{-59}$). It is below $9/32$ for every $N \ge 10^8$: on $\ell \in [\ell_0, 4\cdot10^4]$, where $\ell_0 = \log(10^8/\pi) = 17.2759\ldots$, the function $T$ is strictly decreasing (an Arb proof over $8{,}000$ boxes, \path{rhi_monotone}), and for $\ell \ge 4\cdot10^4$ one has $T \le 10^{-4} + \log(1 + k/2)/k \le 6.09\cdot10^{-4}$. The same fact, by a cruder route, is the kernel-checked statement of \cref{sec:v2lean}. At $\ell = 10^3$ the bound is $0.0122908$, at $\ell = 10^4$ it is $0.00182589$, and asymptotically it is $(\sqrt5 + o(1))\,\log\log N/\log N$, which is much larger than $1/\log N$, consistent with the Siegel sanity check of \cref{sec:v2checks}. The formula \eqref{eq:C2bound} itself falls below $9/32$ near $N = 10^{7.83}$, but the theorem is stated, and certified, only for $N \ge 10^8$. For comparison, \cref{thm:main} gives $9/32 = 0.28125$ with OpenAI's theorem; \cref{thm:C2} gives $0.27643$ at $N = 10^8$ without it, and less from there on.

\paragraph{How the proof goes (\cite[P5 memo \S A.13.3]{dq1records}).} With $p = \sqrt5$, $q_0 = 1$ and $C = \ell/4 - 1$, so that $k = 4C/\sqrt5 - 8/(5-1) = (4/\sqrt5)(\ell/4 - 1) - 2$ as in \eqref{eq:C2bound}:
\begin{enumerate}
\item On $[0, t_{\min}]$ the maximal zero height contracts classically.
\item On $[t_{\min}, T(\ell)]$ the floor $S_\chi(x,\sqrt5,t) \ge \ell/4 - 1$, for every real $x$ (components 4a--4c below), feeds the one-step comparison \eqref{eq:HBineq}.
\item The first contact of the barrier $\sqrt{Q(t)}$ with a zero is attained, because Lemma Z$'$ (component 5) confines all zeros of height at least $\sqrt b$ to $|x| < X(\ell)$, uniformly in $t$.
\item Hence the maximal squared height is at most $b$ at time $T(\ell)$.
\item De Bruijn's strip theorem gives $\Lc_\chi \le T(\ell) + b/2$.
\end{enumerate}

\paragraph{The components, with their trust classes and records.}\label{sec:v2trust}
\begin{center}\small
\begin{longtable}{@{}P{0.04\linewidth}P{0.30\linewidth}P{0.36\linewidth}P{0.22\linewidth}@{}}
\toprule
\# & Component & Content, certificate & Class; review \\
\midrule
1 & Domain reduction & The floor is needed only for $t \in [t_{\min}(\ell), T(\ell)]$; in $(R,\tau)$ this is a compact region for $\ell \le 4\cdot10^4$ and the small-$\tau$ region $\tau \le 0.01745$, $R \ge 2.618$ for $\ell > 4\cdot10^4$ \cite[P5 memo \S A.9.0]{dq1records} & \ref{T3}; FGMPS Thm.~2.6 \ref{T4}. REVIEW\_P5\_DELTA \\
2 & Gaussian tilt \eqref{eq:tilt} & The exact identity for $S_\chi$; FGMPS (4.1) & \ref{T3}, \ref{T4}. REVIEW\_P5 \\
3 & Transfer & Planat Thm.~2.1 and Prop.~2.3; Hakimoglu-Brown Thm.~9, one step at $p = \sqrt5$. The bridge numerator $(c^2-1)(c^2-5)$ was certified in Telperion and matched to Planat's Prop.~2.3 by the auditor, who did \emph{not} re-check the Telperion certificate & \ref{T4}; transposition \ref{T3}. REVIEW\_P5 \\
4a & Floor, compact part (G1) & $S_\chi(x,\sqrt5,t) \ge \ell/4 - 1$ for $\ell \in [\ell_0, 4\cdot10^4]$, $t \in [10^{-4}, T(\ell)]$, $0 \le x \le 10^4$, both parities. Certifier \lean{y3\_cert.py} (md5 \lean{5b375c5b}) with \lean{y3\_big.py} (\lean{aa232458}), domain \lean{c2\_domain.py} (\lean{f84d38a1}), driver \lean{g1\_sweep\_y3.py} (\lean{a591107d}); four ledgers, $67{,}679$ certified leaves, $0$ capped, $0$ failed, worst certified value $-0.9999990867 \ge -1$; $0.43$ CPU-hours & \ref{T2}. REVIEW\_Y3\_G1: reproduced, identical in every field except wall time \\
4b & Floor, $|x| \ge 10^4$ (LX) & $S_\chi \ge \ell/4 - 0.77$ for $\ell \le 4\cdot10^4$, $t \in [10^{-4}, T(\ell)]$, $|x| \ge 10^4$. Proved $x$-monotonicity (\S A.9.2) plus \lean{y3\_big2.py} (\lean{aafeba32}) over $(R,\tau)$ boxes: \lean{lx\_cover\_p5patched.out} (\lean{5aa52aef}), $3{,}069$ boxes, $0$ failed, worst $-0.7689010122$ & \ref{T2}; monotonicity \ref{T3} (parts (b)--(g) read, not fully re-derived). REVIEW\_Y3, REVIEW\_Y3\_G1: reproduced \\
4c & Floor, $\ell > 4\cdot10^4$ (STL) & $S_\chi \ge \ell/4 - 0.18$ for $\tau \le 0.0175$, $R \ge 2.618$, every $x$; analytic, with constants in Arb ($Z(2.6179) = 0.250988$ gives $-0.17525$). Contains every domain point with $\ell > 4\cdot10^4$ & \ref{T3}. REVIEW\_P5\_DELTA, REVIEW\_P5\_A11 \\
4d & Domain tiling & The upper edge $R_{\mathrm{hi}}(\tau)$ of the compact region is monotone on each side of $\tau^*$: \lean{rhi\_monotone.py} (\lean{66a25bab}), $8{,}000$ boxes, $0$ failures; outward edges & \ref{T2}. REVIEW\_P5\_A11 \\
5 & Confinement (Lemma Z$'$) & $H_t(x+iy) \ne 0$ for $|x| \ge 4e^{8\max(4\cdot10^4,\ell) - \ell}$, $\min(0.01, 1/\ell) \le |y| \le 1$, $t_{\min} \le t \le \tfrac12$. From FGMPS Thm.~4.3 and Prop.~4.18, which FGMPS state for $3 \le N \le 4\cdot10^5$; the extension to every $N$ is P5's reading of their proof, checked at source by the auditor. Constants in Arb: \lean{p2z\_arb.py} (\lean{b732f0ef}), $61{,}440$ boxes, $0$ failed, worst $Q = 0.908 > 0$ & \ref{T3}, \ref{T4}, \ref{T2}. REVIEW\_P5\_A12 \\
6 & Special functions & DLMF 5.15.1, 5.7.6, 5.6.6; BMOR Lemma~2.1, the $|\log L| \le \log\zeta$ step of Prop.~2.2, Lemma~5.6; FGMPS Lemma~4.2(5) (uniform Stirling; also proved directly in \cref{sec:defects}); Hakimoglu-Brown p.~11 & \ref{T4}. REVIEW\_P5 \\
7 & Values, and ``$< 9/32$'' & Arb; and \cref{sec:v2lean} & \ref{T2}; \ref{T1}. REVIEW\_P5\_A12, REVIEW\_LEAN\_DQ1 \\
\bottomrule
\end{longtable}
\end{center}

\noindent
The four G1 ledgers are:
\begin{center}\small
\begin{tabular}{@{}lrrrrrrl@{}}
\toprule
ledger & initial boxes & evaluations & leaves & capped & failed & worst & md5 (DONE) \\
\midrule
\lean{y3g1\_corner\_a0} & 320 & 5,538 & 2,929 & 0 & 0 & $-0.9999656033$ & \lean{79194715} \\
\lean{y3g1\_corner\_a1} & 320 & 378 & 349 & 0 & 0 & $-0.9988738033$ & \lean{0c5a24e4} \\
\lean{y3g1\_rem\_a0} & 4,576 & 74,495 & 37,065 & 0 & 0 & $-0.9999990867$ & \lean{43e1f9bc} \\
\lean{y3g1\_rem\_a1} & 4,576 & 50,165 & 27,336 & 0 & 0 & $-0.9999363845$ & \lean{f630953b} \\
\bottomrule
\end{tabular}
\end{center}
The ``corner'' jobs cover the hard corner $\ell$ near $\ell_0$, small $t$, $x \le 5$, where the margin is thinnest; the ``rem'' jobs cover the rest, and drop boxes that fall outside the domain or inside the STL region. The jobs ran with six workers, box caps of $300$ and $240$ seconds, bisection depth at most $30$, and took $12.6$, $1.0$, $152.9$ and $93.1$ seconds of wall time: about $0.43$ CPU-hours in all. The auditor re-ran all four from copies and obtained the same ledgers in every field except wall time \cite[REVIEW\_Y3\_G1]{dq1records}.

\subsection{Four defects in the certifiers, and the re-run}\label{sec:defects}

We record this because it bears on how much to trust class \ref{T2}. The first version of \cref{thm:C2} was certified by the P5 session's own certifiers, with $440{,}993$ leaves in $2.9$ CPU-hours, and the auditor passed them in four reviews (REVIEW\_P5, \_DELTA, \_A11, \_A12), checking each majorant cell by cell and re-running a $184$-box sample. The next session (Y3), which copied the certifiers for \cref{thm:Y3}, found three defects in them, and the auditor, adjudicating those, found a fourth \cite[Y3 memo \S4; REVIEW\_Y3]{dq1records}:
\begin{enumerate}[label=B\arabic*.]
\item In the G1 certifier, a Riemann sum over a box could extend to points where the real part of the argument of $-\zeta'/\zeta$ fell below $1$; there the ``majorant'' is negative and the enclosure is not established. This \emph{formally invalidated} the G1 certificate. (In the LX certifier the same cells produce a NaN and the box is rejected, so LX failed safe.)
\item In the LX certifier, a lower bound for the Gaussian mass of the bulk was taken at the wrong end of the box, so it was an upper bound; the denominator was overestimated and a box value could come out too high. This \emph{invalidated the LX certificate as run}; the corrected cover passes with the same constant, the worst value moving from $-0.7689010032$ to $-0.7689010038$ with this fix alone.
\item The closed-form majorant for the far tail of LX is not integrable as $U \to -\infty$ (its logarithm grows like $+(1-U)^2/(4v_0)$), so the step ``$\int_{-\infty}^{U_{\min}} \le 2q(U_{\min})/\mathrm{rate}$'' had no proof as written. A proof-step gap with no effect on any value; repaired by a three-region tail argument whose new ingredient, $\log|\Gamma(x+iy)| \ge \log\Gamma(x) - \pi|y|/2 - \tfrac12\log(1+4y^2)$ for $x \ge \tfrac12$, is proved from \cite[5.7.6]{DLMF}. The same region also uses the upper bound $\log\Gamma(y) \le y\log y$ for real $y = (s_1 + a)/2 \ge 1250$, which the records take from FGMPS's uniform Stirling estimate \cite[Lemma 4.2(5)]{FGMPS2026}, read at source: for $|\im z| \ge 1$ or $\re z \ge 1$, $\Gamma(z) = \sqrt{2\pi}\exp\bigl((z-\tfrac12)\log z - z + O_{\le}(1/(12(\max\{\re z, |\im z|\} - \tfrac12)))\bigr)$, which on the real axis gives $\log\Gamma(y) \le y\log y$ for $y \ge 2$. Since that inequality is load-bearing, here is a self-contained proof for every real $y \ge 1$ (PROVED; it is not in the records). Write $y = n + \theta$ with $n = \lfloor y\rfloor \ge 1$ and $0 \le \theta < 1$. Then $\Gamma(y) = \Gamma(1+\theta)\prod_{j=1}^{n-1}(\theta + j)$. Each factor $\theta + j \le y$, and $\Gamma(1+\theta) \le 1$ because $\log\Gamma$ is convex (H\"older's inequality in the integral $\int_0^\infty t^{y-1}e^{-t}\,dt$) with $\Gamma(1) = \Gamma(2) = 1$. So $\Gamma(y) \le y^{\,n-1} \le y^{\,y-1}$, and $\log\Gamma(y) \le (y-1)\log y \le y\log y$. (For $y < 1$ the inequality fails, since $\log\Gamma(y) > 0 > y\log y$ there; the range $y \ge 1250$ is what region IIb uses.)
\item (Found by the auditor.) In LX, after the contour shift, the bulk Riemann sum started at a point that is an \emph{upper} bound for the true start of the bulk, so a set of points was in neither the sum nor the tail. A coverage gap, proof step only; repaired by a uniform bound for the missing set, written by the auditor, re-derived independently by Y3, and adopted. With it the LX worst value is $-0.7689010122$.
\end{enumerate}
All four repairs were judged sound by the auditor, who is not an independent check of B4 and says so. With the repaired G1 certifier the whole G1 cover was re-run, with the record's jobs and parameters, and passed with $67{,}679$ leaves instead of $440{,}993$: the B1 clamp removes a near-pole majorant that had forced deep bisection. In $1{,}087$ measured box pairs the patch never lowered a floor and sometimes raised one. The P5 outputs are superseded and kept as history (\path{provenance/P5_SUPERSEDED.md}); the certificates of record are those in the table above, and \cref{thm:C2} rests on them. The auditor's own account of what it had failed to check \cite[REVIEW\_Y3 \S7]{dq1records} is four items: that every majorant is used inside its domain at every point of a ball; that every bound feeding a denominator has the right direction over the whole ball; that a closed-form tail is integrable and not merely of the right local slope; and that a bulk/tail split is the same set on both sides when a contour shift has an unknown real part. We state them here because they are the kind of error that ball arithmetic does not catch: the arithmetic was rigorous, and the enclosure it was asked to compute was the wrong one.

\subsection{The arithmetic step, kernel-checked}\label{sec:v2lean}

The single statement of version~2 that is a Lean theorem is that the bound of \cref{thm:C2} is below $9/32$ for every $N \ge 10^8$. The file \lean{DQ1\_P5.lean} (md5 \lean{01a24ec6}, $346$ lines; \path{lean/v2/} in the release) defines, in a namespace \lean{C2}, the heat flow in its Gaussian form \eqref{eq:gaussform} on top of Mathlib's \lean{DirichletCharacter.LFunction} and \lean{Complex.Gamma},
\begin{lstlisting}
def xi (χ : DirichletCharacter ℂ m) (s : ℂ) : ℂ :=
  ((m : ℂ) / (Real.pi : ℂ)) ^ ((s + (parity χ : ℂ)) / 2) * Complex.Gamma ((s + (parity χ : ℂ)) / 2) *
    DirichletCharacter.LFunction χ s
def H (χ : DirichletCharacter ℂ m) (t : ℝ) (z : ℂ) : ℂ :=
  ∫ v : ℝ, xi χ ((1 + I * z) / 2 + ((Real.sqrt t * v : ℝ) : ℂ)) * ((Real.exp (-v ^ 2) / Real.sqrt Real.pi : ℝ) : ℂ)
\end{lstlisting}
together with the functions of \eqref{eq:C2bound}, transcribed symbol for symbol,
\begin{lstlisting}
def Lc (m : ℕ) : ℝ := Real.log ((m : ℝ) / Real.pi)
def tmin (L : ℝ) : ℝ := min (1 / 10000) (4 / L)
def bC2 (L : ℝ) : ℝ := min (1 / 10000) (1 / L ^ 2)
def kC2 (L : ℝ) : ℝ := 4 / Real.sqrt 5 * (L / 4 - 1) - 2
def TC2 (L : ℝ) : ℝ :=
  tmin L + Real.log ((2 + kC2 L * (1 - 2 * tmin L)) / (2 + kC2 L * bC2 L)) / kC2 L
def C2bound (m : ℕ) : ℝ := TC2 (Lc m) + bC2 (Lc m) / 2
\end{lstlisting}
and proves
\begin{lstlisting}
theorem C2bound_lt_9_32 (hm : 10 ^ 8 ≤ m) : C2bound m < 9 / 32
\end{lstlisting}
by pure real analysis, with no Arb input: $T \le 10^{-4} + f(k)$ with $f(k) = \log(1 + k/2)/k$, $f$ is antitone on $(0,\infty)$, $\ell \ge \ell_0 \ge 17.25$, hence $k \ge 3.925$ and $b \le 17.25^{-2}$, and $10^{-4} + 1.0862/3.925 + 1/(2\cdot17.25^2) = 0.27852 < 0.28125$. This is cruder than Arb's $0.27643$ but it is what the theorem needs. The file also proves \lean{STL\_applies} (that every domain point with $\ell > 4\cdot10^4$ satisfies the hypotheses of the STL lemma), and \lean{H\_conj}, \lean{S\_neg} and \lean{TL\_floor\_C2\_of\_nonneg} (the reflection $S_\chi(-x) = S_\chi(x)$ and the reduction of the floor to $x \ge 0$), the last three under the explicit hypothesis $L(\bar s,\chi) = \overline{L(s,\chi)}$, which is true for real $\chi$ but is not in Mathlib. The theorem of \cref{thm:C2} itself, the floor, the confinement lemma, LX and STL are recorded in the same file as \lean{Prop}-valued definitions, for instance
\begin{lstlisting}
def C2_large_conductor_statement : Prop :=
  ∀ {m : ℕ} [NeZero m], 10 ^ 8 ≤ m → ∀ (χ : DirichletCharacter ℂ m), χ.IsPrimitive → IsRealChar χ →
    ∀ t : ℝ, C2bound m ≤ t → ∀ z : ℂ, H χ t z = 0 → z.im = 0
\end{lstlisting}
with no \lean{axiom} and no \lean{sorry}: they are stated, not proved, and nothing in the file claims them.

\emph{What this check consists of.} The file was built once by the DQ-1 Lean builder session at Lean \lean{v4.35.0} with a Mathlib of the same version, in a private environment that is \emph{not} the OpenAI environment of \cref{sec:formal}: exit code $0$, no errors, and \lean{\#print axioms} reporting \lean{[propext, Classical.choice, Quot.sound]} for each of the five proved declarations. The auditor reviewed the file for faithfulness to the statement of record, symbol by symbol, and checked the arithmetic of the proof route, but did not run Lean (REVIEW\_LEAN\_DQ1). There was no five-checker replay, no Comparator run and no tamper test. There is also no formal connection between \lean{C2.H} and version~1's \lean{DBN.Gen.Hchi}: the former is the Gaussian form \eqref{eq:gaussform} on Mathlib's \lean{LFunction} and is defined for $t > 0$ (Mathlib's $\Gamma$ takes a junk value at its poles, which lie on a null set of the integration variable), the latter is the cosine transform of the theta kernel; that they are the same flow up to the factor $8$ is FGMPS's (4.1) and \cref{sec:normalisation}, by hand. So the kernel-checked content of version~2 is the unconditional arithmetic fact that $T(\ell) + b/2 < 9/32$ for every $N \ge 10^8$; that this quantity bounds $\Lc_\chi$ is \cref{thm:C2}, in classes \ref{T2}--\ref{T4}. The middle range of \cref{sec:W2B3} has no Lean counterpart.

\subsection{Theorem W2B3: the middle range, extending Theorem C2 downwards}\label{sec:W2B3}

\emph{This subsection is new in version 2.1.0.} \Cref{thm:C2} starts at $N = 10^8$ because one contraction step from the trivial strip needs $k \ge 3.765$, and Remark~\ref{rem:U2} shows how far one step with a lower probe height reaches ($N \ge 15{,}848{,}932$ for odd and $N \ge 36{,}307{,}806$ for even $\chi$). The result below moves the unconditional wall down by a further decade. It was obtained on 2--3 November 2026 by one agent session (``W2B3''), whose memo is the record cited below, and reviewed by the auditor (REVIEW\_W2B3 \cite{dq1records}), who found one defect in the new certifier, required its repair and the re-issue of the certificates, and confirmed the re-issued record. \emph{No part of it is kernel-checked.}

\begin{theorem}[W2B3; classes \ref{T2}--\ref{T4}; not kernel-checked]\label{thm:W2B3}
Let $\chi$ be a real primitive Dirichlet character of conductor $N$. If $\chi$ is odd and $1{,}650{,}000 \le N \le 10^8$, then every zero of $H_t$ is real for every $t \ge 0.2807226$. If $\chi$ is even and $3{,}700{,}000 \le N \le 10^8$, then every zero of $H_t$ is real for every $t \ge 0.2808812$. With \cref{thm:C2} for $N \ge 10^8$: $\Lc_\chi < 9/32$ for every odd real primitive $\chi$ with $N \ge 1{,}650{,}000$ and every even one with $N \ge 3{,}700{,}000$.
\end{theorem}

\noindent
\emph{Hypotheses.} Exactly those of \cref{thm:C2}: $\chi$ real and primitive, nothing else. GRH is not assumed, no zero-free half-plane is assumed (the initial strip is the trivial $|\im z| \le 1$), no numerical verification of zeros is used, and a Siegel zero, if there is one, is allowed. The margins to $9/32$ are $5.27\cdot10^{-4}$ (odd) and $3.69\cdot10^{-4}$ (even), lower ends of Arb balls. The constants are the upper ends of $128$-bit balls, rounded up; the auditor's independent landing in \lean{mpmath} at $50$ digits gives $0.2807225$ and $0.2808811$, agreeing to $10^{-7}$ (REVIEW\_W2B3). Against Remark~\ref{rem:U2} the thresholds fall by factors $9.6$ and $9.8$.

\paragraph{The derivative kernel (cited).} The pairwise inequality behind \eqref{eq:HBineq} compares the kernel $K_h$ of a zero pair at the height $h$ of the highest zero with the same kernel at the probe height $p$. Hakimoglu-Brown's Lemma~10 \cite[p.~14, (25)--(26)]{HakimogluBrown2026}, read at source by the W2B3 session and by the auditor, adds a derivative term. Let $p > 0$, $B = p^2$, $0 < A^2 < B$ and $\lambda, \mu \ge 0$; for $D \ge 0$ and $0 \le r \le A^2$ put $w = D + B - r$, $\Delta = w^2 + 4Dr$, $\mathcal N = 2\lambda w\Delta - \mu(2w + 4B)\Delta + 8\mu Bw^2$ and
\begin{equation}\label{eq:HBphi}
\Phi(D, r) = 2\Delta^2 - (D + A^2 + 3r)\,\mathcal N .
\end{equation}
If $\Phi \ge 0$ on that domain, then for $0 < h^2 \le A^2$, $0 \le v \le h$ and $(d, v) \ne (0, h)$,
\[
K_h(d, v) \;\ge\; h\Bigl(\frac{\lambda}{p}K_p(d, v) - \mu\,\partial_pK_p(d, v)\Bigr),\qquad K_p(d, v) = \frac{p - v}{d^2 + (p-v)^2} + \frac{p + v}{d^2 + (p+v)^2},
\]
with $D = d^2$ and $r = v^2$ (so $K_p = 2pw/\Delta$). At $(\lambda, \mu) = (1, 0)$ this is the inequality used for \cref{thm:C2}. Summing over the zero pairs, $S_\chi(x, p, t) = \sum K_p$ and $\partial_pS_\chi = \sum\partial_pK_p$, so a lower bound $C$ for $S_\chi$ \emph{and an upper bound $G$ for $\partial_pS_\chi$} on a time slab give, by his Theorem~12 \cite[p.~17, (33)]{HakimogluBrown2026}, the contraction $q' \le -2 - kq$ for the squared height $q$ of the highest zero, with
\begin{equation}\label{eq:k33}
k \;=\; 4\Bigl(\frac{\lambda C}{p} - \mu G - \frac{2\lambda}{B - q_i} - \frac{2\mu(B + q_i)}{(B - q_i)^2}\Bigr) \;>\; 0
\end{equation}
on a slab on which $q \le q_i \le A^2$; the last two terms are the removed own pair, and the bracket decreases in $q$ (his p.~17). His Theorem~12 is a stepped comparison: slabs $t_0 < t_1 < \dots$, constants $(C_i, G_i, \lambda_i, \mu_i)$ and a certified $\Phi_i \ge 0$ per slab, and the barrier recursion $q_{i+1} = (q_i + 2/k_i)e^{-k_id_i} - 2/k_i$ with $d_i = t_{i+1} - t_i$, until the barrier reaches $b$ inside a slab at the landing time of \eqref{eq:landing}. In his paper the slabs end at $t \le 1/5$, $b = 1/25$, and the floors hold only beyond a finite region; here, as in \cref{thm:C2}, the floor and the ceiling are uniform in $x$, the slabs run to $t = 0.29$, $b = 10^{-4}$, and $(\lambda_i, \mu_i)$ vary with the slab, which his proof allows because it is an induction over slabs. Everything else in the assembly is as in \cref{thm:C2}: the classical contraction on $[0, t_{\min}]$ from $q_0 = 1$, the compact attainment of the first contact by Lemma~Z$'$ for $3 \le N \le 10^8$ (\lean{p2z\_small}, \cref{sec:Y3}), and multiple zeros and zeros on $i\R$ through FGMPS's Proposition~3.2 and Lemma~3.7; real zeros enter with half-pair weight, as in his proof, and Lemma~10 is per pair and linear with non-negative weights, so it does not interact with that splitting. The transposition is class \ref{T3}; the auditor checked Lemma~10, Theorem~12 and \eqref{eq:k33} at source and re-derived $K_p$, $\partial_pK_p$, the bracket and the own-pair terms.

\paragraph{The new quantity, and its certifier.} From the tilt identity \eqref{eq:tilt}, since the logarithm of the prefactor is linear in $c$ and $dc/dz = -i/2$,
\[
\partial_pS_\chi(x, p, t) \;=\; \tfrac14\,\re\Bigl(\frac{D'}{D}\Bigr)'(c) .
\]
Write $f = \Phi_a'/\Phi_a$ and $\langle\,\cdot\,\rangle$ for the normalised Gaussian average that the certifiers of \cref{sec:v2method} already enclose, as a bulk sum plus tails, so that $D'/D = \langle f\rangle$ and $D''/D = \langle f' + f^2\rangle$. For any constant $m$,
\[
(D'/D)' \;=\; \langle f' + (f - m)^2\rangle - \langle f - m\rangle^2 ,
\]
and with $m$ the value of $f$ at the centre of the average (for LX, the centre of the double contour shift) the two averages are small: on the bulk the mean value theorem on the segment from the centre gives $|f(w) - m| \le \tau|g|\sup|\psi_a'| + \min(\tau|g|\sup|(L'/L)'|,\ |L'/L(w)| + |L'/L(c)|)$ and $|f'(w)| \le \sup|\psi_a'| + \sup|(L'/L)'|$, with $(\zeta'/\zeta)'$ as the majorant of $|(L'/L)'|$ (decreasing on $(1, \infty)$) and $\psi'$ for $|\psi_a'|$ (decreasing, \cite[5.15.1]{DLMF}, evaluated rigorously with \lean{acb\_polygamma}); the tails of $\Phi_a''$ are integrated by parts twice, with the same three far-tail regions as the repair B3 of \cref{sec:defects}. The result is $G \le \tfrac14(E_1 + E_2^2)$ with $E_1$, $E_2$ upper bounds for the two averages. The two certifiers, \lean{w2b3\_cert.py} (md5 \lean{f2e56a85}) for $0 \le x \le 10^4$ and \lean{w2b3\_big2.py} (\lean{4b780036}) for $x \ge 10^4$, extend \lean{y3\_cert.py} and \lean{y3\_big2.py} by accumulators only: the floor they return for $S_\chi$ is bit-identical to the version-2.0.0 certifiers' on $16$ G1 boxes and $10$ LX boxes (\lean{w2b3\_check.py}; re-run for this release). Two parameters were changed for both quantities: the bulk/tail split moved from $\delta = 1/4$ to $\delta = 1/2$, which is still inside the range $-\delta \le \sigma \le 1 + \delta$, $\delta \le 1/2$ of Rademacher's bound \cite[Lemma 5.6]{BMOR2020} and is asserted by the sweep; and, for the derivative terms only, the sharper $\Gamma$-majorant $\psi'$ in place of the trigonometric one. One new monotonicity claim is needed for LX, that its boundary term at the split is non-increasing in $x$ for $x \ge 10^4$, which holds when $R - u_s \ge 1/2$; at $p = 7/2$ one has $R \ge 9/4$ and $u_s = 3/2$, and the certifier asserts it. The auditor scanned the claim numerically on $[2500, 2.5\cdot10^{15}]$ and read the derivation. As a calibration, not a proof, the certified $G$ lies above the exact value at $x = 0$, $\tfrac14(\zeta'/\zeta)'(R) + \tfrac1{16}\psi'(\tfrac{R + a}{2})$, everywhere sampled, by $0.03$ at $t \approx 0$ and by up to $0.11$ at larger $t$.

\paragraph{The covers.} Probe height $p = 7/2$ ($B = 49/4$), $\delta = 1/2$, and $39$ contiguous time slabs per parity with edges $10^{-4}, 0.005, 0.010, \dots, 0.100, 0.11, \dots, 0.29$, each certified over the band $[N_{\min}, 10^8]$ and every $x \ge 0$ (negative $x$ by the reflection, which holds for $\partial_pS_\chi$ as for $S_\chi$), with the $\tau$-edges padded outward so that the certified $t$-ranges overlap at every seam. A box passes when $S_\chi - \ell/4 \ge \mathrm{CF}_i$ \emph{and} $G \le \mathrm{CG}_i$; the driver \lean{w2b3\_sweep.py} (\lean{9efa5464}) is \lean{y3\_sweep.py}'s with the second test added, bisecting to depth $24$ (G1) and $16$ (LX). The targets $\mathrm{CF}_i$, $\mathrm{CG}_i$ are the slab extremes measured on a fixed grid (\lean{w2b3\_measure.py}, $136$ fine slabs) minus and plus $0.004$, and the per-slab $(\lambda_i, \mu_i)$, multiples of $1/50$, were chosen by \lean{w2b3\_design.py} to maximise $k_i$ subject to the exact certificate of $\Phi_i$ below, with $A_i^2 = \lceil 1000q_i\rceil/1000 + 1/1000 \ge q_i$. The design thresholds were $N = 1.587\cdot10^6$ and $3.599\cdot10^6$; the certified bands start above them for margin.
\begin{center}\small
\begin{tabular}{@{}llrrrrl@{}}
\toprule
parity & band & slabs & boxes (evaluations) & failed & wall & md5 \\
\midrule
odd & $[1{,}650{,}000, 10^8]$ & 39 & 4,926 (6,078) & 0 & 126 s & \lean{696148f1} \\
even & $[3{,}700{,}000, 10^8]$ & 39 & 5,925 (8,076) & 0 & 166 s & \lean{0be5d4cf} \\
\bottomrule
\end{tabular}
\end{center}
The md5 is that of the $39$ \lean{.DONE} files of the parity concatenated in slab order, and ``boxes'' counts the certified leaves. Each ledger records \lean{complete: true}, \lean{n\_failed: 0}, the exact floats $\mathrm{CF}_i$, $\mathrm{CG}_i$ the boxes were compared against, the $N$-edges rounded outward in $\ell$, the $t$-edges, and the md5 prefixes of the certifier code that produced it. The two covers ran with two workers each, about five CPU-minutes in all; the whole result, measurements and design included, cost about $1.45$ CPU-hours.

\paragraph{The positivity certificates.} Hakimoglu-Brown's (34)--(36) \cite[pp.~22--23]{HakimogluBrown2026} certify $\Phi_i \ge 0$ exactly: $\Psi_i(U, V) = ((1 - U)^4/B^4)\,\Phi_i(BU/(1-U), A_i^2V)$ is a polynomial of degree at most $4$ in each variable on $[0, 1]^2$; on each of eight strips in $U$ its $25$ Bernstein coefficients are computed, and if all are positive then $\Psi_i > 0$ on the square, hence $\Phi_i > 0$ for every finite $D \ge 0$ and $0 \le r \le A_i^2$, with no separate large-$D$ argument (the corner $U = 1$ is the leading coefficient $2 - 2\lambda + 2\mu$). \lean{w2b3\_phi.py} (\lean{797ff1d8}) does this in Python's exact rationals, with no floating point. The certificates of record, \lean{w2b3\_phi\_a1\_m1650000.json} (\lean{4e72c655}) and \lean{w2b3\_phi\_a0\_m3700000.json} (\lean{c3213386}), have $39$ rows each, all certified, smallest coefficient $1.34\cdot10^{-4}$. Self-tests: row $0$ of his Table~3 ($p = 61/20$, $q_0 = 9/16$, $\lambda = 1.718$, $\mu = 0.868$) certifies with smallest coefficient $7.1\cdot10^{-3} > 13/10000$, as his (36) states; $(\lambda, \mu) = (1, 0)$ at $B = 49/4$, $A^2 = 1$ gives exactly $0$, the boundary case of his remark after Lemma~10; and the negative control $(5/2, 0)$ is refused (smallest coefficient $-3$). The auditor's own exact implementation reproduced the smallest coefficient of every one of the $78$ rows exactly. Two rows of a first design, checked only on a float grid, were refused by the exact certificate ($A^2 = 23/1000$ and $3/1000$); the plans of record were designed under the exact certificate, and the landing refuses any uncertified row.

\paragraph{The landing.} \lean{w2b3\_bound.py} (\lean{4b4ca675}), in Arb at $128$ bits with outward rounding: $q_0 = 1 - 2\cdot10^{-4}$ at $t_{\min} = 10^{-4}$, $b = 10^{-4}$; per slab $C_i = \ell_{\min}/4 + \mathrm{CF}_i$, $G_i = \mathrm{CG}_i$, $k_i$ by \eqref{eq:k33} with the certificate's rationals; asserted: every ledger complete with no failure and with $N_{\min}$ inside its band, $0 < q_i \le A_i^2 < B$, $k_i > 0$, $t_0 \le 10^{-4}$, and overlapping seams. The barrier reaches $b$ in slab $38$ ($t \in [0.28, 0.29]$) for both parities, at $T = 0.2806725$ (odd) and $0.2808311$ (even), and $\Lc_\chi \le T + b/2$. Since $k_i$ increases with $\ell$, the value at $N_{\min}$ bounds the whole band, as in \cref{sec:Y3}. Outputs: \lean{w2b3\_bound2\_a1.json} (\lean{87e26183}) and \lean{w2b3\_bound2\_a0.json} (\lean{55c37c93}); re-run from the released ledgers for this version, byte-identical. Every fourth slab:
\begin{center}\small
\begin{tabular}{@{}lrrrrrrrrr@{}}
\toprule
& $i$ & $t_i$ & $\mathrm{CF}_i$ & $\mathrm{CG}_i$ & $A_i^2$ & $\lambda_i$ & $\mu_i$ & $q_i$ (up) & $k_i$ (down) \\
\midrule
odd & 0 & 0.0001 & $-0.20265$ & 0.21709 & 1001/1000 & 8/5 & 16/25 & 0.99980 & 3.4212 \\
& 8 & 0.0400 & $-0.17840$ & 0.19966 & 199/250 & 81/50 & 33/50 & 0.79463 & 3.5829 \\
& 16 & 0.0800 & $-0.15345$ & 0.18947 & 307/500 & 83/50 & 18/25 & 0.61246 & 3.7221 \\
& 24 & 0.1400 & $-0.11652$ & 0.17538 & 191/500 & 44/25 & 9/10 & 0.38018 & 3.9260 \\
& 32 & 0.2200 & $-0.06485$ & 0.14887 & 7/50 & 9/5 & 24/25 & 0.13855 & 4.2313 \\
& 38 & 0.2801 & $-0.03162$ & 0.12956 & 3/1000 & 46/25 & 26/25 & 0.00133 & 4.4378 \\
\midrule
even & 0 & 0.0001 & $-0.34252$ & 0.25610 & 1001/1000 & 8/5 & 16/25 & 0.99980 & 3.4348 \\
& 8 & 0.0400 & $-0.32132$ & 0.24103 & 199/250 & 81/50 & 33/50 & 0.79449 & 3.5830 \\
& 16 & 0.0800 & $-0.29351$ & 0.23314 & 307/500 & 83/50 & 18/25 & 0.61246 & 3.7137 \\
& 24 & 0.1400 & $-0.25109$ & 0.22046 & 191/500 & 42/25 & 37/50 & 0.38043 & 3.9123 \\
& 32 & 0.2200 & $-0.18923$ & 0.19170 & 141/1000 & 9/5 & 24/25 & 0.13896 & 4.2262 \\
& 38 & 0.2801 & $-0.14840$ & 0.16875 & 3/1000 & 46/25 & 26/25 & 0.00165 & 4.4537 \\
\bottomrule
\end{tabular}
\end{center}
The plain kernel $(\lambda, \mu) = (1, 0)$, fed the same ledgers, gives $k_i$ between $2.82$ and $3.17$ and does not land at all before $t = 0.29$ (the barrier is still at $q \approx 0.037$ and $0.032$ there): every decade gained over Remark~\ref{rem:U2} is the derivative kernel's.

\paragraph{The review, and a fifth defect.} The auditor read both certifier diffs in full and the relevant pages of Hakimoglu-Brown at source, and ran its own exact Bernstein checker, its own landing, a monotonicity scan and, from a copy of the certifier, all $78$ slab sweeps. It found one defect in the new G1 accumulator (``F1''): the mean value theorem on the segment from the centre to a bulk point needs the supremum of $|(L'/L)'|$ over the segment, that is, the majorant at the segment's smallest real part, and the first code took it at the point's own real part, which, the majorant being decreasing, under-estimates for points to the right of the centre. The defect was not load-bearing: with the one-line fix all $78$ slabs were re-swept as the record ($6{,}078$ and $8{,}076$ evaluations), with no failure, the largest rise in a slab's certified $G$ being $2.6\cdot10^{-3}$ (odd) and $3.1\cdot10^{-3}$ (even); the targets $\mathrm{CF}_i$, $\mathrm{CG}_i$, the slab edges and the positivity certificates are unchanged, so the landing files are byte-identical to the pre-fix ones. The pre-fix code and ledgers are superseded and are not in the release. The auditor's own re-run of the fixed code agrees with the re-issued ledgers in every field except wall time, in all $78$ slabs (REVIEW\_W2B3, Addendum~1). The fix consumed most of the $0.004$ slack on $G$: the smallest remaining $\mathrm{CG}_i - \max G$ is $1.9\cdot10^{-5}$ (odd, slab $8$) and $3.3\cdot10^{-5}$ (even, slab $7$). That is thin, but it is a certificate, because $\mathrm{CG}_i$ is used in the landing as the exact binary float the boxes were compared against. We record F1 beside B1--B4 of \cref{sec:defects}: it is the same kind of error, a bound used at the wrong point of a ball, which the ball arithmetic cannot see.

\paragraph{What this theorem is not.} It is not kernel-checked, and unlike \cref{thm:C2} it has no Lean statement at all; the kernel-checked arithmetic fact of \cref{sec:v2lean} is about the explicit function of \cref{thm:C2} and covers only $N \ge 10^8$. Lemma~Z$'$ for $N > 4\cdot10^5$ rests on our reading of FGMPS's Theorem~4.3 and Proposition~4.18 beyond their stated range, as for \cref{thm:C2}. Hakimoglu-Brown's Lemma~10 and Theorem~12 are in an unrefereed working draft about $\zeta$; their use here is an AI-written transposition reviewed by one AI reviewer. The certifier programs are the trust base of class \ref{T2}, and the review found a fifth defect in them. Nothing here bears on GRH.

\subsection{Theorem Y3: every conductor, given the half-plane}\label{sec:Y3}

\begin{theorem}[Y3; classes \ref{T2}--\ref{T4}; conditional on OpenAI's half-plane, as \cref{thm:main}]\label{thm:Y3}
Let $\chi$ be a real primitive Dirichlet character of conductor $N \ge 3$, and assume that $L(s,\chi) \ne 0$ for $\re s > 7/8$. Then every zero of $H_t$ is real for every $t \ge 0.27912$. More precisely, for $N < 10^8$,
\[
\Lc_\chi \le 0.2791114 \ (\chi \text{ odd}),\qquad \Lc_\chi \le 0.2783189 \ (\chi \text{ even}),
\]
the worst cases being $N = 3$ and $N = 5$, and for $N \ge 10^8$ \cref{thm:C2} gives $\Lc_\chi \le 0.2764282$ with no hypothesis at all.
\end{theorem}

\noindent
This is a strict improvement on version~1's $9/32 = 0.28125$ for every conductor, by at least $2.1\cdot10^{-3}$. The hypothesis is discharged exactly as in \cref{thm:main}, by OpenAI's kernel-checked theorem; with the functional equation it gives the strip $|\im z| \le 3/4$, that is, $q_0 = 9/16$, and $t_{\min} = b = 10^{-4}$. The proof is the method of \cref{sec:v2method} with the probe height $p = 20$ and the floors below, from the Y3 session \cite[Y3 memo \S3]{dq1records}; the auditor re-ran the two certificates and obtained ledgers identical in every field except wall time, recomputed the landing independently in \lean{mpmath} at $50$ digits, and checked the sources named above at source (REVIEW\_Y3 \S6).

\begin{center}\small
\begin{tabular}{@{}lllllllr@{}}
\toprule
ledger & parity & $N \in$ & $t \in$ & floor $S_\chi - \ell/4 \ge$ & worst G1 / LX & leaves & md5 \\
\midrule
\lean{u20\_a1} & odd & $[3, 10^8]$ & $[10^{-4}, 0.499]$ & $0.385$ & $0.3903$ / $1.811$ & 3,672 & \lean{4a1c199c} \\
\lean{u20\_a0} & even & $[5, 10^8]$ & $[10^{-4}, 0.499]$ & $0.36$ & $0.3640$ / $1.811$ & 3,672 & \lean{953e869c} \\
\bottomrule
\end{tabular}
\end{center}
Each took about $40$ seconds, with $0$ failures. (At $N = 3$ one has $\ell = \log(3/\pi) < 0$, which the certifier's domain handling accommodates; the auditor checked this.) The landing \eqref{eq:landing} with $k = 4(\ell/4 + \mathrm{CF})/p - 8/(p^2 - q_0)$ is decreasing in $k$, and $k$ increases with $\ell$, so each band's worst case is its lowest conductor. The auditor's independent landing values are:
\begin{center}\small
\begin{tabular}{@{}llllll@{}}
\toprule
case & $N$ & $k$ & $T$ & $\Lc_\chi \le$ & margin \\
\midrule
odd, $q_0 = 9/16$ & 3 & 0.054666 & 0.2790613 & 0.27911134 & $9/32 - \Lc_\chi \ge 2.14\cdot10^{-3}$ \\
even, $q_0 = 9/16$ & 5 & 0.075207 & 0.2782689 & 0.27831886 & $2.93\cdot10^{-3}$ \\
odd, $q_0 = 1$ & 3 & 0.054644 & 0.4932441 & 0.49329408 & $\tfrac12 - \Lc_\chi \ge 6.71\cdot10^{-3}$ \\
even, $q_0 = 1$ & 5 & 0.075185 & 0.4907846 & 0.49083455 & $9.17\cdot10^{-3}$ \\
C2 at $N = 10^8$ & $10^8$ & 3.937186 & 0.2763782 & 0.27642820 & $4.82\cdot10^{-3}$ \\
\bottomrule
\end{tabular}
\end{center}
Two labels, which we keep as the records keep them: the even constant $0.2783189$ is the \emph{certificate value}, computed from the exact binary floating-point floor $0.36$ that the sweep compared against, rounded up. The records also carry \emph{Lean-rational} values, from the rational floors $3599/10000 \le 0.36$ and $385/1000 \le 0.385$ chosen for a future Lean statement: exactly $0.27831962506\ldots$ (even) and $0.27911134136\ldots$ (odd), so, rounded up, $0.2783197$ and $0.2791114$. (The records display them to seven digits rounded to nearest, as $0.2783196$ and $0.2791113$; those displays are not upper bounds. The Lean statements sketched in the records use $27833/100000$ and $27912/100000$, which are safe.) The theorem above uses the certificate values. The confinement lemma for this range, Lemma Z$'$ for $3 \le N \le 10^8$, was certified separately (\lean{p2z\_small}: $76{,}800$ boxes, $0$ failed, worst $Q \ge 0.9455$); for $N \le 4\cdot10^5$ FGMPS's Theorem~2.7 also gives it directly.

\emph{Why a height of $20$.} At height $p$ the $\Gamma$ factor alone contributes about $\tfrac14\psi(\tfrac{1+p}{4} + \tfrac a2) \approx \tfrac14\log(p/4)$ to $S_\chi - \ell/4$, and this beats the loss $8/(p^2 - q_0) \approx 8/p^2$ in $k$ once $p$ is large, for \emph{every} $N \ge 3$; the arithmetic enters only through the small term $\tfrac12\zeta'/\zeta(\tfrac{1+p}{2})$. The price is that $k$ is small ($0.055$ to $0.075$ at the worst conductors), so the gain over $9/32$ is small, of the order of $kq_0^2/8$. The Y3 memo also records five certified bands with smaller $p$ in the middle range, which give smaller constants there (for instance at most $0.2424$ at $N = 10^4$ and $0.2139$ at $N = 10^6$ for either parity, \cite[Y3 memo \S3.4]{dq1records}); we do not restate them. We also do not use the per-character minimum of these bounds with FGMPS's unconditional ones, which would lower the uniform constant to $0.27832$, because it would rest on FGMPS's Table~10, which we have read but not reproduced.

\subsection{Theorem U1: every conductor, unconditionally}\label{sec:U1}

\begin{theorem}[U1; classes \ref{T2}--\ref{T4}; no hypothesis]\label{thm:U1}
Let $\chi$ be a real primitive Dirichlet character of conductor $N \ge 3$. Then every zero of $H_t$ is real for every $t \ge 0.49330$: $\Lc_\chi \le 0.4932941$ for odd $\chi$ and $\Lc_\chi \le 0.4908346$ for even $\chi$ when $N < 10^8$ (certificate values, rounded up), and $\Lc_\chi \le 0.2764282$ when $N \ge 10^8$.
\end{theorem}

\noindent
This is the same pair of certificates as \cref{thm:Y3}, with the trivial initial strip $q_0 = 1$ in place of $9/16$ (the certificate's $t$-range $[10^{-4}, 0.499]$ contains the landing time $0.49324$, and $p = 20 \ge \sqrt5$). The gain over de Bruijn's $1/2$ is $6.7\cdot10^{-3}$, which is small; the point is what the statement does not use.

\emph{What is new here, checked at source.} FGMPS write, in their survey of previous work: ``On the upper-bound side, only the trivial bound $\Lc_L \le \tfrac12$ from de Bruijn's argument was known. Theorem 1.1 gives the first non-trivial upper bounds for a Dirichlet $L$-function'' \cite[\S1]{FGMPS2026}. Their Theorem~1.1 is stated for $3 \le N \le 400$ and uses Platt's verification of GRH to a height $h_N \ge 10^8/N$; their large-$x$ results (Theorems~2.7, 4.3, 8.1) are stated for $N \le 4\cdot10^5$ and give zero localisation, not a bound on $\Lc_L$. So FGMPS leave every $N > 400$ at $1/2$ without further hypotheses. With OpenAI's half-plane granted, version~1 of this paper (and FGMPS's Remark~3.4) already give $9/32 < 1/2$ for every $N$. The accurate claim is therefore: \cref{thm:U1} is the first bound below $1/2$ that is uniform in the conductor and uses no zero-free half-plane, no GRH and no numerical zero verification. A search of the literature (alphaXiv, by the Y3 session) found no other upper bound on $\Lc_L$ for Dirichlet $L$-functions; that is a search result, not a proof of absence.

\begin{remark}[C2's threshold lowered]\label{rem:U2}
The same machinery at $p = 5/2$, with floors certified from measured worst values ($-0.55$ for odd and $-0.76$ for even $\chi$; ledgers \lean{c25\_a1}, \lean{c25\_a0}, reproduced by the auditor), gives unconditionally $\Lc_\chi \le 0.2811138 < 9/32$ for odd $\chi$ with $N \ge 15{,}848{,}932$ and $\Lc_\chi \le 0.2812404 < 9/32$ for even $\chi$ with $N \ge 36{,}307{,}806$ \cite[Y3 memo \S3.5]{dq1records}. The margins are $1.36\cdot10^{-4}$ and $9.67\cdot10^{-6}$: thin, but decided, since the floor is used as the exact binary float the sweep compared against. The one-step method from $q_0 = 1$ cannot go below about $N = 10^{7.03}$ (odd) or $10^{7.34}$ (even) even with the exact $t = 0$ floor; the rest of the gap is certifier slack. (Version 2.1.0: \cref{thm:W2B3} lowers these thresholds to $1{,}650{,}000$ and $3{,}700{,}000$ with thirty-nine steps and the derivative kernel; the statement, ledgers and constants of this remark stand as published.)
\end{remark}

\subsection{Consistency checks}\label{sec:v2checks}

None of these is part of a proof.
\begin{itemize}
\item \emph{Against FGMPS.} For the seven smallest conductors, FGMPS's unconditional bounds (their Table~10, $D = 1$) and \cref{thm:Y3}'s conditional ones are: $\chi_{-3}$: $0.26967$ vs $0.27911$; $\chi_{-4}$: $0.26$ vs $0.27856$; $\chi_5$: $0.27912$ vs $0.27832$; $\chi_{-7}$: $0.28097$ vs $0.27748$; $\chi_{-8}$: $0.26007$ vs $0.27723$; $\chi_8$: $0.26343$ vs $0.27742$; $\chi_{-11}$: $0.28605$ vs $0.27663$. Both are upper bounds, so no contradiction is possible; the two methods are complementary, and for the $93$ characters with $N \le 400$ for which FGMPS's $D = 3/4$ search keeps the strip bound $9/32$ \cite[\S7.7]{FGMPS2026}, \cref{thm:Y3} improves all of them.
\item \emph{Against direct computation.} For six actual characters ($\chi_{-3}, \chi_{-4}, \chi_5, \chi_{-7}, \chi_8, \chi_{-11}$) the Y3 session computed $S_\chi(x, 20, t)$ directly in \lean{mpmath} from \eqref{eq:tilt} on a grid of $t$ and $x$; the minimum of $S_\chi - \ell/4$ was $0.030$ above the certified floor everywhere sampled, and equal to the exact $t = 0$ floor to $10^{-3}$. So the certificate is tight, and the $p = 20$ route is at its structural limit.
\item \emph{Siegel sanity.} Every bound here is a constant at least $0.18$, and no rate is $o(1/\log N)$; a method that gave $\Lc_\chi = o(1/\log N)$ uniformly without excluding Siegel zeros would be suspect. The method also gives nothing for the ``zoo'' of fake $L$-functions with the right functional equation but off-line zeros, because they violate its Euler-product inputs; the floor itself never sees zeros.
\end{itemize}

\subsection{Reproducing the certificates}\label{sec:v2repro}

The \path{certificates/} directory of the release holds the as-run sources, drivers, ledgers and outputs, each with its md5 in \path{certificates/MANIFEST.md}, and a README with the exact commands. The certifiers need Python~3 with \lean{python-flint} and run one core per worker, with the BLAS pinned to one thread. Costs measured on the record runs: the four G1 jobs, $0.43$ CPU-hours with six workers; the LX cover, $36$ seconds; each of the four Y3 band ledgers, $40$ to $70$ seconds; the two confinement checks, $2$ seconds each; the Lean file, $18$ seconds at Lean \lean{v4.35.0}. For \cref{thm:W2B3} (\path{certificates/w2b3/}): the $78$ slab sweeps, $126$ and $166$ seconds of wall time with two workers each; the two positivity certificates and the landing, under a second each; the bit-identity check, one second; the two measurement runs behind the targets, about ten minutes each. A reader who re-runs a certifier should compare the ledger's \lean{evaluations}, \lean{certified}, \lean{worst} and \lean{failed} fields (for \path{w2b3/}, also \lean{maxG\_g1} and \lean{maxG\_lx}) with the record's; the auditor's reproductions matched every field except wall time.

%======================================================================
\section{Limitations}\label{sec:limits}
%======================================================================

\begin{enumerate}
\item \emph{OpenAI's manuscript is unrefereed.} The $7/8$ half-plane is a kernel-checked Lean theorem, but its mathematical proof has not been refereed, and FGMPS treat it as ``asserted''. Our theorem is exactly as reliable as that Lean proof and the definitions it is stated in.
\item \emph{The replays are local and unsandboxed.} All the checks of \cref{sec:verification} ran on one macOS machine, outside any sandbox, driven by AI agent sessions. The sandboxed Comparator run (\cref{sec:comparator}) closes this gap only in part: it ran as uid 0 with all capabilities dropped, Mathlib came from its official build cache, \lean{lean4lean} could not take part, and it was made by another agent session: the auditor read its logs and reproduced the solution export byte for byte, but did not repeat its kernel runs.
\item \emph{Cross-version checking.} The export was made with the v4.34.1 exporter and checked by checkers bundled with v4.35.0-rc3, which read the same export format. \lean{lean4lean} could not take part.
\item \emph{Mathlib's definitions are trusted.} In particular \lean{IsPrimitive}, \lean{LFunction} and \lean{completedLFunction} are taken to mean what their docstrings say. Part (b) of \cref{thm:main} makes our own definition of $H^\chi_t$ answer to Mathlib's $L$-function, but it cannot check Mathlib.
\item \emph{Gauss's sign theorem is proved for real primitive characters only.} That is all the proof needs. We do not prove anything about the root number of a complex character, and the remark in \cref{sec:gauss} about avoiding the sign altogether is unformalized.
\item \emph{The constant $\Lc_\chi$ is not defined in Lean} (Remark~\ref{rem:sinf}). Nor is the identification of $H^\chi_t$ with FGMPS's $H^L_t$ for $t \ne 0$, which is by hand (\cref{sec:normalisation}).
\item \emph{The local build tree's dependencies were not hashed.} See \path{provenance/OAI_ENV.md} in the release. The kernel replays check every proof, but the definitions the statement uses were taken as they stood in that tree. For the main theorem's closure this gap is closed, relative to Mathlib's official build cache at \lean{d13f23b}, by the byte-identical export of \cref{sec:comparator}.
\item \emph{The audits were done by AI sessions.} No human has read the Lean proofs; the kernel checks them, and a human should read the statement and \lean{ChallengeDefs} (\cref{sec:statement}).
\item \emph{The numerics are checks, not proofs.} Nothing in the formal proof depends on them.
\item \emph{Version 2's bounds are not kernel-checked.} \Cref{thm:C2,thm:Y3,thm:U1,thm:W2B3} are in classes \ref{T2}--\ref{T4}: Arb certificates, AI-written proofs reviewed by one AI reviewer, and unrefereed preprints. The one kernel-checked statement of \cref{sec:v2lean} is the arithmetic of the explicit function of \cref{thm:C2}, for $N \ge 10^8$ only, checked by one kernel in one build, with no further checkers and no Comparator run. The Lean statement layer sketched for \cref{thm:Y3,thm:U1} in the Y3 memo was not built, and \cref{thm:W2B3} has no Lean statement.
\item \emph{The certifiers are programs.} Ball arithmetic guarantees the arithmetic, not the enclosure asked for. Four defects (\cref{sec:defects}) were found in certifiers that four reviews had passed, and a fifth (\cref{sec:W2B3}) in the derivative certifier of version 2.1.0; the repaired certifiers were reviewed and their outputs reproduced, but another defect of the same kind would not be caught by re-running them. A reader who wants to trust class \ref{T2} should read the certifier sources.
\item \emph{The contraction inequalities are preprints about $\zeta$.} Planat's Theorem~2.1 and Proposition~2.3, and Hakimoglu-Brown's Theorem~9 and paired-zero inequality, and, for \cref{thm:W2B3}, his Lemma~10 and Theorem~12, are stated for Riemann's $\xi$; Hakimoglu-Brown's paper is a working draft. Their use for $H_t$ of a Dirichlet character, with a different confinement lemma and different constants, is an AI-written transposition. FGMPS's Theorem~4.3 and Proposition~4.18 are stated for $3 \le N \le 4\cdot10^5$; their use for every $N$ is our reading of their proof, checked at source by the auditor but by no one else.
\item \emph{Two steps were read, not re-derived, by the auditor:} the Telperion-certified bridge numerator in component~3, and parts (b)--(g) of the $x$-monotonicity lemma behind LX (component~4b). The auditor says so in its reviews, and we repeat it here.
\item \emph{No formal link between the two heat flows.} \lean{C2.H} (version~2) and \lean{DBN.Gen.Hchi} (version~1) are different Lean definitions, in different environments; that they define the same flow up to the factor $8$ is by hand (\cref{sec:normalisation}).
\item \emph{Some margins are thin.} Remark~\ref{rem:U2}'s even case is below $9/32$ by $9.67\cdot10^{-6}$, and in \cref{thm:W2B3} the certified ceilings on $\partial_pS_\chi$ exceed the largest certified value by as little as $1.9\cdot10^{-5}$ in one slab. The decisions rest on using the floor or ceiling as the exact binary float the certifier compared against; the Lean-rational constants are chosen on the safe side of those floats.
\item \emph{Nothing here is GRH progress}, in either version. The unconditional bounds are consistent with the existence of Siegel zeros, and the method does not see zeros at all.
\end{enumerate}

%======================================================================
\section{Use of AI}\label{sec:ai}
%======================================================================

This work was done with AI agent sessions (Claude, by Anthropic) under the direction of the author. We describe who did what as precisely as we can.

\begin{itemize}
\item \emph{Proof engineering.} An agent session (the ``contrarian'') wrote the Lean sources of the \lean{Gen} layer, the Gauss layer and \lean{GaussBridge}, and the generator that derives four \lean{Gen} files from the island. The \lean{dbn} island itself was written in September and October 2026 with AI assistance under the author's direction, as recorded in \cite{Murphy932,MurphyNewman}.
\item \emph{Audit and composition.} A second, separate agent session (the auditor, ``bb'') reviewed each layer as it arrived. It read the new sources in full and rebuilt the whole chain from source (OpenAI's library was used as built in the 7 October replay, with \path{Nonvanishing.lean} re-elaborated), and wrote its own check files, tampered statements and closure walks. It wrote the numerical checks of \cref{sec:verification}, wrote the two composition files and ran all the checker replays. Its reviews are labelled ``AI review, NOT an operator read-back''. They are the records cited as \cite{bbrecords}. It also wrote the Comparator challenge files of \cref{sec:comparator} and read that run's seven logs in full, matched their md5s, and reproduced its solution export byte for byte.
\item \emph{Coordination and this paper.} A third session (the ``lead'') coordinated the work and wrote the publication plan. This paper was drafted by a further agent session, from the auditor's records and from sources it read. Every number in it is traceable to those records or sources, and anything it could not verify is marked as unresolved.
\item \emph{Version 2.} The large-conductor bound (\cref{thm:C2}), the Gaussian-tilt identity, the domain reduction, the three certifiers and the first certificates were the work of one agent session (``P5''), in the course of a research dive on a different question, whose memo records the statement of record (\S A.13). A second session (``Y3'') obtained \cref{thm:Y3,thm:U1}, found defects B1--B3, wrote the repairs, and ran the G1 re-run. The auditor (``bb'', the same auditor session lineage as in version~1) reviewed P5 in four rounds, adjudicated B1--B3, found B4 and wrote its repair, reproduced every certificate of record (identical in every field except wall time), reviewed the Lean file for faithfulness, and reviewed the whole package before release. A Lean-builder session wrote and built \lean{DQ1\_P5.lean}. The lead session ruled that the release form is a new version of this record rather than a separate note. This version of the paper was drafted by a further agent session from those records and from the sources, each read at source; every number in \cref{sec:v2} is traceable to a ledger or a source, and what could not be verified is marked.
\item \emph{Version 2.1.0.} One agent session (``W2B3'') read Hakimoglu-Brown's Lemma~10, Proposition~11 and Theorem~12 and his Bernstein certificate at source, wrote the derivative certifiers as extensions of the version-2.0.0 certifiers, the measurement, design, positivity-certificate and landing programs, ran the covers, and wrote the memo. The auditor (``bb'', the same lineage) read the certifier diffs in full and the cited pages at source, wrote its own exact positivity checker and its own landing, re-ran all $78$ slab sweeps from a copy, found the defect F1, required the repair and the re-issue of the record, and confirmed the re-issued ledgers field by field. This version of the paper was revised by a further agent session from those two records and the sources, which re-ran the landing, the positivity certificates and the bit-identity check from the released files and re-swept every slab from the released code; every number in \cref{sec:W2B3} is traceable to a ledger or a source.
\item \emph{The author's role.} The author directed the programme: the author decided which results to pursue, approved running the work in OpenAI's environment and the publication of this result, chose the venue, and read and approved the final text. For version~2 the author approved the release of the large-conductor bound and the rest of the lead's list on 1 November 2026, and read and approved this text on 1 November 2026. For version 2.1.0 the author approved the release of the middle range on 3 November 2026, and read and approved this text on 3 November 2026.
\end{itemize}

\noindent
The Lean kernel is the trusted component. No proof in this work is accepted on the word of an AI session. According to their READMEs, two of the checkers used in \cref{sec:verification} were themselves written with AI under human supervision (Joachim Breitner, Lean FRO): \lean{con-leche}, including its consistency proof, by Claude, and its Rust port \lean{con-ron} by an AI system the README does not name \cite{conleche,conron}; the five acceptances of \cref{sec:verification} also include Lean's own kernel and \lean{nanoda}, which do not share their code. The definitions and the statement are what a human reader must check, and \cref{sec:statement} is written to make that check feasible. The definitions and the statement (\cref{sec:statement}) were read by the author before publication; no outside reader has checked them.

%======================================================================
\section*{Acknowledgements}
%======================================================================

The bound is Freire, Gonz\'alez-Riquelme, Mu\~noz, Poblete and Salinas's conditional observation; the engine is de Bruijn's theorem; the half-plane and its Lean proof are OpenAI's; the sign is Gauss's. Mathlib's $L$-function library made the bridge to $\Lc(s,\chi)$ a statement about a shared definition rather than a private one. The independent checker \lean{nanoda} is Chris Bailey's; \lean{lean4export} and Comparator are from the Lean FRO (Comparator by Henrik B\"oving). The checker \lean{con-leche} \cite{conleche} was written by Claude, and its Rust port \lean{con-ron} \cite{conron} by an AI system its README does not name, both under the supervision of Joachim Breitner at the Lean FRO. Cureton and Thorner's half-plane and Lean development make the companion section possible. Version~2 adds debts: the contraction inequalities are Planat's and Hakimoglu-Brown's; the Gaussian form of the heat flow, the strip theorem in its time-shifted form, the finiteness theorem and the effective approximation are FGMPS's; the Euler-product and convexity bounds for $L(s,\chi)$ were read in Bennett, Martin, O'Bryant and Rechnitzer, who quote Rademacher; the digamma and Gamma inequalities are from the DLMF; the ball arithmetic is Fredrik Johansson's Arb, through \lean{python-flint}. Version 2.1.0 adds one more: the conjugate-pair inequality with a derivative term, the stepped derivative-kernel comparison and the Bernstein form of its positivity certificate are Hakimoglu-Brown's.

\bibliographystyle{amsplain}
\bibliography{refs}

\appendix
\section{The auditor's records}\label{app:records}

The records cited as \cite{bbrecords} are the following files. They are archived, each with its md5, in the \path{provenance/} directory of the release repository:
\begin{itemize}\small
\item \lean{REVIEW\_\allowbreak{}GEN\_\allowbreak{}D1\_\allowbreak{}D2a.md}: the generic de Bruijn library, and the theta kernel;
\item \lean{REVIEW\_\allowbreak{}D2b\_\allowbreak{}BRIDGE.md}: the Mellin bridge;
\item \lean{REVIEW\_\allowbreak{}D2c\_\allowbreak{}STRIP.md}: the zero strip;
\item \lean{REVIEW\_\allowbreak{}D2bp\_\allowbreak{}CHI4.md}: the theta functional equation, and $\chi_{-4}$;
\item \lean{NOTE\_\allowbreak{}D2\_\allowbreak{}GAUSS\_\allowbreak{}SIGN.md}: Mathlib's root number, and the sign;
\item \lean{REPORT\_\allowbreak{}OAI\_\allowbreak{}PORT.md}: the port to OpenAI's environment, and provenance;
\item \lean{REVIEW\_\allowbreak{}GAUSS\_\allowbreak{}SIGN\_\allowbreak{}AND\_\allowbreak{}OAI.md}: Gauss's sign theorem, and the composition;
\item \lean{REVIEW\_\allowbreak{}GAUSS\_\allowbreak{}BRIDGE\_\allowbreak{}ADD1.md} (``ADD1''): part (b), and the conjunction;
\item \lean{REVIEW\_\allowbreak{}P1P6\_\allowbreak{}BUNDLE.md}: including the literature status of the statement (its \S6);
\item \lean{REVIEW\_\allowbreak{}CTQ.md}, \lean{REVIEW\_\allowbreak{}INGHAM.md}, \lean{REVIEW\_\allowbreak{}CTFREE\_\allowbreak{}KERNELS.md}: the companion results;
\item \lean{REVIEW\_\allowbreak{}CT\_\allowbreak{}RERUNS.md}: the capability-dropped Comparator re-runs of P2, Ingham and CT main;
\item \lean{COMPARATOR\_\allowbreak{}RUN\_\allowbreak{}SPEC.md}: the specification of the sandboxed run;
\item \lean{REVIEW\_\allowbreak{}COMPARATOR\_\allowbreak{}LAMBDA\_\allowbreak{}CHI.md}: the verification of that run, with \lean{REVIEW\_\allowbreak{}COMPARATOR\_\allowbreak{}LAMBDA\_\allowbreak{}CHI\_\allowbreak{}ADDENDUM1.md}, the inspection of its logs, and \lean{REVIEW\_\allowbreak{}COMPARATOR\_\allowbreak{}LAMBDA\_\allowbreak{}CHI\_\allowbreak{}ADDENDUM2.md}, two precision corrections.
\end{itemize}

\section{The records of version 2}\label{app:v2records}

The records cited as \cite{dq1records} are the following files, archived with their md5s in \path{provenance/} of the release (the first eight characters of each md5 are given):
\begin{itemize}\small
\item \lean{P5\_\allowbreak{}memo\_\allowbreak{}extract.md} (\lean{4b3688f2}): sections A.3 and A.6--A.13 of the P5 session's memo, reproduced verbatim; the full memo (md5 \lean{e3676c4c}, $1{,}775$ lines) is held internally, since most of it concerns other questions, and the extract's header records the line ranges taken. \S A.3 and \S A.6, the $t = 0$ floor and the Gaussian tilt; \S A.7, the explicit floor; \S A.9, the domain reduction, the double contour shift, LX and STL; \S A.12.2, Lemma Z$'$; \S A.13, the statement of record of \cref{thm:C2} with its component table and the Lean statement layer. The reviews cite the full memo by its md5 and by section number, and the section numbers are unchanged in the extract;
\item \lean{P5\_\allowbreak{}SUPERSEDED.md} (\lean{f723042e}): which P5 outputs are superseded by the re-runs, and by what;
\item \lean{Y3\_\allowbreak{}memo\_\allowbreak{}extract.md} (\lean{ebb8f66b}): sections 0--5, 8 and 9 of the Y3 session's memo, reproduced verbatim; the full memo (md5 \lean{3779ce49}, $720$ lines) is held internally, and the two omitted sections (\S6--\S7) are internal dive bookkeeping that nothing here cites. \S2, where C2 fails below $10^8$ and why; \S3, \cref{thm:Y3,thm:U1} and Remark~\ref{rem:U2}; \S4, defects B1--B4, their repairs, and the G1 re-run (\S4.6); \S5, the consistency checks; \S8, a Lean statement layer, not built; \S9, the file and md5 table;
\item \lean{REVIEW\_\allowbreak{}P5.md} (\lean{e909a59c}), \lean{REVIEW\_\allowbreak{}P5\_\allowbreak{}DELTA.md} (\lean{206b88e7}), \lean{REVIEW\_\allowbreak{}P5\_\allowbreak{}A11.md} (\lean{e570fc33}), \lean{REVIEW\_\allowbreak{}P5\_\allowbreak{}A12.md} (\lean{f6cade54}), \lean{REVIEW\_\allowbreak{}P5\_\allowbreak{}FINAL.md} (\lean{9ae4f16c}): the auditor's four reviews of P5 and its delta on the statement of record;
\item \lean{REVIEW\_\allowbreak{}Y3.md} (\lean{799a8fa8}): the adjudication of B1--B3, the finding of B4, the gate for the G1 re-run, and the review of \cref{thm:Y3,thm:U1} with their reproductions;
\item \lean{REVIEW\_\allowbreak{}Y3\_\allowbreak{}G1.md} (\lean{4b4d8aab}): the confirmation of the G1 re-run and of the B4 adoption, with reproductions of all four ledgers and of the LX cover, identical in every field except wall time;
\item \lean{REVIEW\_\allowbreak{}LEAN\_\allowbreak{}DQ1.md} (\lean{e4e14877}): the faithfulness review of \lean{DQ1\_P5.lean}, among other files of the same build;
\item \lean{REVIEW\_\allowbreak{}PRERELEASE\_\allowbreak{}v2.0.0.md} (\lean{d3a2e9eb}): the auditor's pre-release review of the version-2.0.0 package, which also re-ran \lean{p2z\_arb} and \lean{rhi\_monotone} (components 4d and 5) from copies and obtained byte-identical outputs;
\item version 2.1.0: \lean{W2B3\_\allowbreak{}memo.md} (\lean{c5cd8f39}, $293$ lines): the W2B3 session's memo, with the erratum F1 in its \S0 and the re-issued record in its \S5, \S6 and \S9; and \lean{REVIEW\_\allowbreak{}W2B3.md} (\lean{e195ebaa}): the auditor's review (PASS after the required fix F1), with its Addendum~1 confirming the re-issued record and its Addendum~2 correcting its own figure for the even-parity rise in $G$ to $3.07\cdot10^{-3}$. Both records are in \path{provenance/}, unredacted, as the others are.
\end{itemize}
The certificates themselves (sources, drivers, ledgers, outputs) are in \path{certificates/}, with every md5 in \path{certificates/MANIFEST.md}; those of version 2.1.0 are in \path{certificates/w2b3/}, whose README says what each file is and how to re-run the cheap checks. Paths of the form \path{~/rh-team/...} inside these records refer to the programme's internal working tree at the time of writing; they are not part of the release, and the copies are not redacted, so that their md5s match the originals.

\end{document}
